% Options for packages loaded elsewhere
\PassOptionsToPackage{unicode}{hyperref}
\PassOptionsToPackage{hyphens}{url}
\PassOptionsToPackage{dvipsnames,svgnames,x11names}{xcolor}
%
\documentclass[
  10pt,
]{article}
\usepackage{amsmath,amssymb}
\usepackage{iftex}
\ifPDFTeX
  \usepackage[T1]{fontenc}
  \usepackage[utf8]{inputenc}
  \usepackage{textcomp} % provide euro and other symbols
\else % if luatex or xetex
  \usepackage{unicode-math} % this also loads fontspec
  \defaultfontfeatures{Scale=MatchLowercase}
  \defaultfontfeatures[\rmfamily]{Ligatures=TeX,Scale=1}
\fi
\usepackage{lmodern}
\ifPDFTeX\else
  % xetex/luatex font selection
  \setmainfont[]{DejaVu Serif}
  \setsansfont[]{DejaVu Sans}
  \setmonofont[]{DejaVu Sans Mono}
  \setmathfont[]{Latin Modern Math}
\fi
% Use upquote if available, for straight quotes in verbatim environments
\IfFileExists{upquote.sty}{\usepackage{upquote}}{}
\IfFileExists{microtype.sty}{% use microtype if available
  \usepackage[]{microtype}
  \UseMicrotypeSet[protrusion]{basicmath} % disable protrusion for tt fonts
}{}
\makeatletter
\@ifundefined{KOMAClassName}{% if non-KOMA class
  \IfFileExists{parskip.sty}{%
    \usepackage{parskip}
  }{% else
    \setlength{\parindent}{0pt}
    \setlength{\parskip}{6pt plus 2pt minus 1pt}}
}{% if KOMA class
  \KOMAoptions{parskip=half}}
\makeatother
\usepackage{xcolor}
\usepackage[margin=2.0cm]{geometry}
\usepackage{longtable,booktabs,array}
\usepackage{calc} % for calculating minipage widths
% Correct order of tables after \paragraph or \subparagraph
\usepackage{etoolbox}
\makeatletter
\patchcmd\longtable{\par}{\if@noskipsec\mbox{}\fi\par}{}{}
\makeatother
% Allow footnotes in longtable head/foot
\IfFileExists{footnotehyper.sty}{\usepackage{footnotehyper}}{\usepackage{footnote}}
\makesavenoteenv{longtable}
\setlength{\emergencystretch}{3em} % prevent overfull lines
\providecommand{\tightlist}{%
  \setlength{\itemsep}{0pt}\setlength{\parskip}{0pt}}
\setcounter{secnumdepth}{-\maxdimen} % remove section numbering
% doc_header.tex -- LaTeX preamble additions used by build_docs.sh for every
% Markdown -> LaTeX -> PDF conversion (the documents of the folder Revision_Lord_Gravity).
\usepackage{etoolbox}
% longtable may not be loaded by pandoc when a document has no table
\usepackage{longtable}
% bold vectors: \bm -> \symbf under unicode-math (XeLaTeX), bm package otherwise
\makeatletter
\@ifpackageloaded{unicode-math}{\providecommand{\bm}[1]{\symbf{#1}}}{\usepackage{bm}}
\makeatother
% smaller type in long tables so that wide audit tables fit the text width;
% inside a table a line break is allowed after an underscore of plain text
% (pandoc writes it as \_), so that long identifiers wrap inside their column
% rows are set a little apart, so that symbols with indices do not touch
\AtBeginEnvironment{longtable}{\footnotesize\renewcommand{\arraystretch}{1.45}%
  \let\sptunderscore\_\def\_{\sptunderscore\allowbreak}}
\setlength{\LTcapwidth}{\textwidth}
% A short table is kept on one page.  The build (code/scripts/keep_tables.py) ends the rows of such a table with \\*,
% which forbids a page break after a row, and sets \sptkeptpenalty to 10000 in front of the table.  This penalty is
% put here in the place of the breakpoint (penalty 0) that longtable leaves after every table, so that the paragraph
% that follows a kept table, its caption, begins on the same page; it is set back to 0 after the table.
\newcount\sptkeptpenalty
\sptkeptpenalty=0
\makeatletter
\patchcmd{\endlongtable}{\endgraf\penalty\z@\addvspace\LTpost}%
  {\endgraf\penalty\sptkeptpenalty\addvspace\LTpost\global\sptkeptpenalty=\z@}%
  {}{\PackageError{doc_header}{the end of longtable could not be patched}{The tables can still be typeset, but a short table may be separated from its caption.}}
\makeatother
% allow line breaks after slashes and hyphens inside table cells
\usepackage{xurl}
% code blocks: break over-long lines instead of running into the margin.
% Plain ``` blocks are typeset by pandoc as {verbatim}, highlighted ones
% (```bash) as {Highlighting}; both are re-declared as breakable fancyvrb
% environments.  The continuation line is marked with a small hooked arrow.
\usepackage{fvextra}
\RecustomVerbatimEnvironment{verbatim}{Verbatim}{breaklines,breakanywhere,fontsize=\small}
\makeatletter
\@ifundefined{Highlighting}{}{%
  \RecustomVerbatimEnvironment{Highlighting}{Verbatim}{breaklines,breakanywhere,fontsize=\small,commandchars=\\\{\}}}
\makeatother
% inline code (pandoc writes it as \texttt{...}): never hyphenate -- an inserted
% hyphen could be read as part of a file name -- but allow a line break after
% an underscore, slash, comma, colon or equals sign, so that long paths and
% column lists wrap inside the text width and inside table columns.
% Inside a formula (e.g. the name of a notebook symbol in an equation) the same
% command must first leave math mode, which \text of amsmath does.
\ExplSyntaxOn
\tl_new:N \l__spt_code_tl
\cs_new_protected:Npn \__spt_texttt:n #1
  {
    \tl_set:Nn \l__spt_code_tl {#1}
    \regex_replace_all:nnN { ([\/\,\:\=]) } { \1 \c{allowbreak} } \l__spt_code_tl
    \regex_replace_all:nnN { \c{_} } { \c{_} \c{allowbreak} } \l__spt_code_tl
    \group_begin: \ttfamily \hyphenchar \font = -1 \scan_stop: \l__spt_code_tl \group_end:
  }
\RenewDocumentCommand \texttt { m }
  {
    \mode_if_math:TF { \text { \__spt_texttt:n {#1} } } { \__spt_texttt:n {#1} }
  }
\ExplSyntaxOff
% A figure is written in the Markdown source as an image paragraph followed by
% a caption paragraph ("**Figure S1.** ..."), so that the caption is also
% visible on GitHub.  An image that starts a paragraph is centered.
\usepackage{graphicx}% pandoc loads it only for documents that contain an image
\NewCommandCopy\sptincludegraphics\includegraphics
\RenewDocumentCommand\includegraphics{O{} m}{%
  \ifvmode\begin{center}\sptincludegraphics[#1]{#2}\end{center}%
  \else\sptincludegraphics[#1]{#2}\fi}
% keep figures where they are referenced
\usepackage{float}
\floatplacement{figure}{H}
% running header with the document title
\usepackage{fancyhdr}
\pagestyle{fancy}
\fancyhf{}
\fancyhead[L]{\small PROJECT spacetimes, Revision\_Lord\_Gravity}
\fancyhead[R]{\small\thepage}
\renewcommand{\headrulewidth}{0.4pt}
% 16x16 matrices: tighter columns so that a full matrix fits the text width
\setlength{\arraycolsep}{3.2pt}
% The documents of this project contain long tagged symbols, file paths and
% identifiers that cannot be hyphenated.  A paragraph is therefore allowed to
% have looser lines before TeX gives up and lets a line stick out of the margin.
\AtBeginDocument{\tolerance=2000 \emergencystretch=5em \hbadness=2000}
\ifLuaTeX
  \usepackage{selnolig}  % disable illegal ligatures
\fi
\IfFileExists{bookmark.sty}{\usepackage{bookmark}}{\usepackage{hyperref}}
\IfFileExists{xurl.sty}{\usepackage{xurl}}{} % add URL line breaks if available
\urlstyle{same}
\hypersetup{
  pdftitle={The task of Revision\_Lord\_Gravity: the gravitational coupling of the brochon by E. A. Lord's gauge theory (task C23), rules, decisions, and plan},
  pdfauthor={Prepared for Patrick Lee Nash, Ph.D.},
  colorlinks=true,
  linkcolor={blue},
  filecolor={Maroon},
  citecolor={blue},
  urlcolor={blue},
  pdfcreator={LaTeX via pandoc}}

\title{The task of Revision\_Lord\_Gravity: the gravitational coupling
of the brochon by E. A. Lord's gauge theory (task C23), rules,
decisions, and plan}
\usepackage{etoolbox}
\makeatletter
\providecommand{\subtitle}[1]{% add subtitle to \maketitle
  \apptocmd{\@title}{\par {\large #1 \par}}{}{}
}
\makeatother
\subtitle{Provenance document \texttt{P0} of PROJECT spacetimes, folder
Revision\_Lord\_Gravity}
\author{Prepared for Patrick Lee Nash, Ph.D.}
\date{November 22, 2026}

\begin{document}
\maketitle

\hypertarget{what-this-document-is}{%
\section{1. What this document is}\label{what-this-document-is}}

This document records the task of the folder
\texttt{Revision\_Lord\_Gravity} of PROJECT spacetimes: the author's
request, in my words and with the times at which he made it; the rules
under which the work is done; the decisions of interpretation, each with
its reason; the plan; and the log of the work, of its errors and of
their corrections. It is written so that it can be read alone: every
symbol it uses is defined in it (Section 8).

\hypertarget{the-task}{%
\section{2. The task}\label{the-task}}

\hypertarget{the-hard-rules}{%
\subsection{2.1 The hard rules}\label{the-hard-rules}}

The work of PROJECT spacetimes is done under ten hard rules of the
author. In substance:

\begin{enumerate}
\def\labelenumi{\arabic{enumi}.}
\tightlist
\item
  \textbf{Quality.} The work must be professional, competent, clear,
  correct, complete and well organized. The readers are first-year
  physics Ph.D.~students who do not know the solution; they must be able
  to learn it from the documents and to teach it to others. Manuscripts
  are always accompanied by provenance.
\item
  \textbf{Honesty and exactness.} ``DO NOT LIE!'' A false statement must
  be explained in a provenance (why it was made, and how it was
  corrected). Instructions must be followed exactly; if one cannot be
  followed, a provenance must explain why. No shortcuts: either the
  required reasoning, creation and calculation are done, or the work
  stops after an explanation in a provenance.
\item
  \textbf{Three versions of every document.} Every documentation, symbol
  and provenance file exists as Markdown (\texttt{.md}), LaTeX
  (\texttt{.tex}) and PDF (\texttt{.pdf}). ``Create a provenance'' and
  ``with provenance'' mean: create these three files.
\item
  \textbf{No referrals.} The reader is never sent to another manuscript,
  document, provenance or symbol file; every equation, symbol and
  definition that a file creates or uses is written in that very file.
  The README, which lists the contents of the folder, is the one
  exception.
\item
  \textbf{Push at once.} Every object is pushed to the repository as
  soon as it is created or updated.
\item
  \textbf{Working directory:}
  \texttt{/Users/nsh/Developer/github/spacetimes}.
\item
  \textbf{Repository:}
  \texttt{https://github.com/once-ere/spacetimes.git}.
\item
  \textbf{One symbol, one object.} A symbol names one object only, in
  the whole project; the symbols of the folder are registered, and their
  relation to the symbols of the other folders of the project is proved
  by a MASTER provenance of the symbols.
\item
  \textbf{Code for everything.} Every equation, result, figure,
  calculation and prediction comes with commented source code,
  instructions for its use, and provenance; every prediction has
  carefully designed graphs. The order of preference of the languages is
  Rust, C, Python 3, Mathematica. The code is executed as a check, with
  provenance.
\item
  The repetition of Rule 4.
\end{enumerate}

Standing instructions of the author that apply here as well: time
integration is done with the integrators of SUNDIALS (the CVODE of his
Rust port, or another applicable SUNDIALS integrator), never with a
hand-written Runge-Kutta method; the work does not stop to wait for him,
but pauses to push and to document; and every qualifying remark about a
result (an approximation, a fit, an accuracy, a limitation) is explained
(what it means, why it holds, why it matters) in the manuscript and in a
provenance of the remarks.

\hypertarget{the-request-task-c23}{%
\subsection{2.2 The request (task C23)}\label{the-request-task-c23}}

On 21 November 2026, at 18:02 PST, the author gave the next task, to be
undertaken after the stages of the revision of the hypotheses then in
progress were complete and pushed. In my words:

\begin{enumerate}
\def\labelenumi{\arabic{enumi}.}
\tightlist
\item
  Read, understand and remember the paper of E. A. Lord, ``General
  relativity from gauge invariance'', Proceedings of the Cambridge
  Philosophical Society 69, 423 (1971).
\item
  Associate the group \(SL(2,\mathbb C)\) of the project with Lord's
  gauging of the representation of a Clifford algebra.
\item
  Introduce a gravitational coupling of the spinor field of the project,
  which is to receive a name, by Lord's covariant derivative (a
  connection added to the partial derivative).
\item
  Analyze, discuss and document this coupling, with provenance.
\item
  Solve the field equations of the spinor in the Schwarzschild field and
  in the Vaidya metric (the metric of a mass that changes by absorbing
  or radiating null dust), and study the behaviour of the solutions on
  both sides of the horizon.
\item
  Make predictions that involve the zeugon and the spinor, such that
  they can be detected gravitationally.
\item
  Push every item as soon as it is complete.
\end{enumerate}

At 18:45 PST he named the path to investigate: the pair of two spinors
bound by the exchange of the zeugon. The binding of such a pair by
photons is, in his words, ``extremely unlikely'' and ``not worth
studying''. He used for the spinor the name \textbf{brochon}, which I
had proposed.

The stages of the revision of the hypotheses were completed in the night
of 21 to 22 November 2026; the last of them was committed and handed to
the author to push on 22 November at 04:35 PST, and this task began
then.

\hypertarget{the-starting-point-the-earlier-hypotheses}{%
\subsection{2.3 The starting point: the earlier
hypotheses}\label{the-starting-point-the-earlier-hypotheses}}

The task is carried out under the hypotheses of the project as they
stood before the revision for the Higgs field of the Standard Model (the
``earlier hypotheses''), because only under them does the zeugon exist.
They are restated here as far as this record needs them. The manuscript
of the folder states them in full since its rewriting of 22 November
2026 (entry E22), in its Appendices A and B, where they are called the
brochon model; before that it did not, and until the correction of
Section 4.1 this sentence said that it did, which was false.

\begin{itemize}
\tightlist
\item
  \textbf{Space-time} is the manifold M31 of the Hermitian
  bi-quaternions, Minkowski space with the metric
  \(\eta^{\rm LL}=\mathrm{diag}(1,-1,-1,-1)\). A unit bi-quaternion is a
  complex \(2\times2\) matrix of determinant 1, so the unit
  bi-quaternions form the group \(SL(2,\mathbb C)\); it is a subgroup of
  the Lorentz super-loop (the unit bi-octonions, a Moufang loop) and
  acts on M31 by the Lorentz transformations.
\item
  \textbf{The brochon} is the spinor field \(\Psi_{\rm LL}\) of the
  earlier hypotheses: an isodoublet of four-component Dirac spinors
  (eight complex anticommuting components) made of the spinor charges of
  the super-loop, with the mass \(m_{\rm LL}\). It carries no electric
  charge and no electroweak quantum numbers.
\item
  \textbf{The gauge group} is
  \(\mathrm{GEE}=SL(2,\mathbb C)\times SU(2)_{\rm v}\times U(1)^{\rm F}_{\rm LL}\):
  the local Lorentz transformations of the spinor index, the isospin
  rotations and the phases of the fermion number. Its covariant
  derivative of the brochon adds to the partial derivative a spin
  connection (the gauge field of \(SL(2,\mathbb C)\)), the three gauge
  fields of \(SU(2)_{\rm v}\) with the coupling \(g^{\rm v}_{\rm LL}\),
  and the gauge field of \(U(1)^{\rm F}_{\rm LL}\) with the coupling
  \(g^{\rm F}_{\rm LL}\). On the flat M31 the spin connection is a pure
  gauge.
\item
  \textbf{The superfluid and the zeugon.} A contact interaction pairs
  two brochons into spin-zero pairs, which form an isospin triplet, and
  the pairs condense. The condensate is polar: it points along one
  isospin direction, which can be taken as the third. It breaks
  \(SU(2)_{\rm v}\times U(1)^{\rm F}_{\rm LL}\) and leaves unbroken the
  rotations about that direction (with a reflection, the group
  \(O(2)\)). The gauge boson of the unbroken generator \(T^3_{\rm LL}\)
  stays massless; it is the \textbf{zeugon} (from the Greek for a yoked
  pair). It couples to the isospin charge \(Q_{\rm LL}=T^3_{\rm LL}\),
  which is \(\pm\tfrac12\) on the two isospin components of the brochon,
  and it attracts the two brochons of a condensed pair, whose charges
  are opposite. It is not the photon.
\end{itemize}

\hypertarget{what-lords-paper-provides}{%
\subsection{2.4 What Lord's paper
provides}\label{what-lords-paper-provides}}

Lord observes that the free field equations of particles of any spin,
written with Dirac matrices in the Weyl (chiral) form, are invariant not
only under Lorentz transformations but also under changes of the
representation of the Pauli algebra by unimodular \(2\times2\) matrices,
a group \(SL(2,\mathbb C)\). He makes that group local, in the manner of
Yang and Mills; covariance then requires auxiliary fields, which he
identifies with the potentials of a spin-2 field, and the resulting
self-interacting spin-2 theory is Einstein's general relativity, the
auxiliary fields becoming the components of the affine connection. The
project's spin connection is a gauge field of the same kind on the flat
M31. The task makes it dynamical by Lord's construction.

\hypertarget{how-the-rules-are-kept-in-this-folder}{%
\section{3. How the rules are kept in this
folder}\label{how-the-rules-are-kept-in-this-folder}}

\begin{itemize}
\tightlist
\item
  The manuscript of the folder is \texttt{lord\_gravity.md} (with
  \texttt{.tex} and \texttt{.pdf}), written part by part; this record
  and the README are written with it. All three versions of every
  document are built by \texttt{code/scripts/build\_docs.sh} from the
  sources in \texttt{docs\_src/}, after
  \texttt{code/scripts/assemble\_docs.py} has inserted the numbers that
  the programs computed. A document that sends its reader to another
  document fails the check
  \texttt{code/scripts/check\_self\_contained.py}.
\item
  Every result is computed by a program in Rust (folder \texttt{rust/}),
  checked by an independent program in Python (folder \texttt{python/})
  and, where exact algebra decides it, by a script in the Wolfram
  Language (folder \texttt{wolfram/}); each writes a report of its
  checks into \texttt{data/}. Time integration uses SUNDIALS.
\item
  Lord's paper is an input that is not in the repository; its SHA-256
  checksum is recorded in \texttt{provenance/inputs\_sha256.txt}, so
  that every statement about it refers to one identified copy.
\item
  Every commit is made in the workspace, carried to the author's
  computer as a git bundle, applied there, and pushed by the author.
\end{itemize}

\hypertarget{records-of-false-statements-rule-2}{%
\section{4. Records of false statements (Rule
2)}\label{records-of-false-statements-rule-2}}

Entry E5 records a clarification that is not a correction of a false
statement: equation (11.1) of the manuscript wrote the spinor connection
with the connection of the vectors in Lord's sign; Part III defines the
spin connection of the tetrad with the opposite (standard) sign, and
(11.1) now states both forms. The false statements found so far are
recorded below.

\hypertarget{objects-of-the-earlier-hypotheses-and-of-the-earlier-work-used-without-being-stated-and-terms-not-defined-22-november-2026}{%
\subsection{4.1 Objects of the earlier hypotheses and of the earlier
work used without being stated, and terms not defined (22 November
2026)}\label{objects-of-the-earlier-hypotheses-and-of-the-earlier-work-used-without-being-stated-and-terms-not-defined-22-november-2026}}

\textbf{What I stated.} (a) Section 2.3 of this record said that the
manuscript of the folder states the earlier hypotheses in full. (b)
Section 3 says that a document that sends its reader to another document
fails the check \texttt{code/scripts/check\_self\_contained.py}, and the
report of that check (\texttt{data/hck\_self\_contained\_checks.txt},
check HCK-H1) said that none of the documents of the folder refers its
reader to another document ``for an equation, a symbol or a
definition''.

\textbf{What was true.} The manuscript \texttt{lord\_gravity.md} (Parts
I to III) uses objects that it takes from the earlier hypotheses of the
project and from the author's unpublished manuscript of 16 November 2026
on the Lorentz super-loop (the ``earlier work''), without stating them;
30 lines of its source contain the words ``earlier hypotheses'' or
``earlier work''. Examples:

\begin{itemize}
\tightlist
\item
  the Lorentz super-loop, the bi-octonions and the Moufang loop: Section
  9.1 says that the unit bi-quaternions form a subgroup of the Lorentz
  super-loop of the bi-octonions and defines none of the three. (A
  \emph{loop} is a set with a product that has a unit element and in
  which every equation \(ax=b\) and \(ya=b\) has exactly one solution; a
  \emph{Moufang loop} is a loop that obeys the Moufang identities, for
  example \(z(x(zy))=((zx)z)y\), a weak form of associativity; the
  \emph{bi-octonions} are the octonions with complex coefficients; the
  \emph{Lorentz super-loop} is the Moufang loop of the bi-octonions of
  complex norm 1, which contains the unit bi-quaternions, the group
  \(SL(2,\mathbb C)\).)
\item
  the methods KS and CJT: Section 15.4 says that the two methods give
  the same ground state, and does not say what they are (two methods of
  the earlier hypotheses for the ground state of the interacting
  brochon: a self-consistent field of the kind of Kohn and Sham, and the
  effective action of Cornwall, Jackiw and Tomboulis);
\item
  the potential \(V^{(2)}_{\rm SF}\) and the recursion of which it is a
  step, the polar ground state, and its numbers
  \(M^{\rm CJT,D}_{\rm SF}=0.1622929325\) and
  \(\lvert D^{\rm CJT}_{\rm SF}\rvert=0.3284861084\), quoted without
  derivation;
\item
  the Bogoliubov-de Gennes equation and the Lagrangian density ``of the
  earlier hypotheses'', the cutoff \(\Lambda_{\rm LL}\) used as the unit
  of the numbers, and the values 0.7 and 0.45 of the couplings.
\end{itemize}

The manuscript also names objects without giving an explicit example of
them (for example the brochon, the zeugon and the polar condensate), and
this record uses terms that it does not define (for example the spinor
charges and the method CJT in Sections 2.3 and 8). The check
\texttt{check\_self\_contained.py} searches only for the explicit
wordings by which a text sends its reader elsewhere (an instruction to
see another document, a statement that something is defined in another
manuscript, a provenance document named as a place to read); it cannot
find a text that uses an object of another document without naming that
document, so its report stated more than it had tested.

\textbf{How it was found.} By my audit of this folder after the author's
message of 22 November 2026, 07:26 PST (entry C3), in which he found
these defects in the article of another folder of the project and stated
that all of my manuscripts have them.

\textbf{Provenance.} The sentence of Section 2.3, the sentence of
Section 3 and the check were written in stage 1, commit \texttt{2d030c1}
(22 November 2026, 04:40 PST); the check was adapted from the folder of
the revision of the hypotheses with its patterns unchanged. The uses of
the earlier hypotheses and of the earlier work came with Parts I to III,
commits \texttt{dced788} (05:03 PST), \texttt{d869b98} (05:08 PST) and
\texttt{b17612c} (06:25 PST).

\textbf{Why I made it.} I have no written record of my reasoning. The
likely cause is that I wrote the manuscript as the continuation of the
earlier documents of the project, for a reader who has them, and treated
the earlier hypotheses as known material to be cited rather than
material to be written down; and that I took the check over from the
other folder without asking what it can detect, and reported its result
without its scope.

\textbf{The correction.} Made in two steps. First (entry E6): the
sentence of Section 2.3 was corrected; the report of HCK-H1 states what
the check searched; and the new check HCK-H2 of the same script lists
every line of the manuscript that contains the words ``earlier
hypotheses'' or ``earlier work'' and fails while there is one. Then,
after the author asked for the manuscripts again (C5), the manuscript
was rewritten (entries E22 to E24): it now contains all the material
that it uses from the earlier hypotheses (Appendix B, the brochon model)
and from the earlier work (Appendix A, the algebra of the Lorentz
super-loop), an explicit example of every object that it names, a word
list (its Section 0.1) and a definition of every term at its first use;
HCK-H2 passes.

\textbf{The pause.} Between the two steps the author had instructed to
continue the calculations; the work continued with them (stages 5 to 7a,
entries E7 to E21).

\hypertarget{local-powers-of-the-decay-rates-read-as-limiting-powers-22-november-2026}{%
\subsection{4.2 Local powers of the decay rates read as limiting powers
(22 November
2026)}\label{local-powers-of-the-decay-rates-read-as-limiting-powers-22-november-2026}}

\textbf{What I stated.} Entry E9 (commit \texttt{b88bfa5}) gave the
local powers of the decay rates of the quasi-bound states, measured
between neighbouring computed masses (6.03 for the states \(1S_{1/2}\)
and \(2S_{1/2}\) and 7.92 for \(2P_{1/2}\), between
\(\mathcal M_{\rm LG}m_{\rm LL}=0.05\) and \(0.1\)), and said that for
the states with the larger powers these masses are not small, so that
their local powers ``are not yet the limiting powers of small masses''.
By naming only those states, it implied that the powers 6.03 and 7.92 of
the others were the limiting powers.

\textbf{What was true.} None of those local powers is a limiting power.
The limiting powers are 5 (\(1S_{1/2}\), \(2S_{1/2}\)), 7
(\(2P_{1/2}\)), 9 (\(2P_{3/2}\), \(3P_{3/2}\)) and 11 (\(3D_{3/2}\)),
that is \(4\ell_{\rm LG}+2S_{\rm LG}+4\) with
\(S_{\rm LG}=j_{\rm LG}-\ell_{\rm LG}\): between
\(\mathcal M_{\rm LG}m_{\rm LL}=0.005\) and \(0.01\) the local powers
are \(1S_{1/2}\): 5.017 (limit 5); \(2S_{1/2}\): 5.017 (limit 5);
\(2P_{1/2}\): 7.014 (limit 7); \(2P_{3/2}\): 9.012 (limit 9);
\(3P_{3/2}\): 9.012 (limit 9); \(3D_{3/2}\): 11.012 (limit 11) (entry
E13), in agreement with the small-mass formula of the literature. At
\(\mathcal M_{\rm LG}m_{\rm LL}=0.05\) the decay rate of the state
\(1S_{1/2}\) is still \(1.179\) times the value of that formula, and the
factor grows with the mass fast enough to raise the local power between
0.05 and 0.1 by about one.

\textbf{How it was found.} By me, on 22 November 2026 at about 10:05
PST, after the commit, when I compared the decay rates with the
published small-mass formula (entry E13).

\textbf{Provenance.} The sentence was written in entry E9 and committed
in \texttt{b88bfa5} (09:52 PST); this record first said ``between 09:40
and 09:58 PST'', a time that ends after the commit (Section 4.3).

\textbf{Why I made it.} I read the local powers of the states whose
decay rates the roots resolved at the smallest computed masses as
limiting powers, without testing their convergence and without comparing
them with the literature before writing.

\textbf{The correction.} Entry E9 now states that the local powers are
not limiting powers and gives the limiting powers; entry E13 records the
analysis that establishes them, with the decay rates computed down to
\(\mathcal M_{\rm LG}m_{\rm LL}=0.005\) by a method that resolves them
in double precision (the charge).

\hypertarget{the-times-of-the-log-22-november-2026}{%
\subsection{4.3 The times of the log (22 November
2026)}\label{the-times-of-the-log-22-november-2026}}

\textbf{What I stated.} The log of Section 7 gives the time of every
entry in Pacific Standard Time. For the entries E1, E3 to E6 and E11 to
E19 the stated interval ends later than the commit in which the entry
first appears, and for E20 and E21 the whole interval, ``12:00 to
12:40'', lies after the commit \texttt{e38f21e} (11:54 PST) that
contains them. Entry C4 called the question that the author answered
``my question of 12:00 PST'', and Section 4.2 gave a time that ends
after its commit.

\textbf{What was true.} The times were estimated from my count of the
work, not read from a clock. The times of the commits (git) are:
\texttt{2d030c1} 04:40 (E1), \texttt{dced788} 05:03 (E2),
\texttt{d869b98} 05:08 (E3), \texttt{b17612c} 06:25 (E4, E5),
\texttt{b88bfa5} 09:52 (E6 to E12), \texttt{0e3ab7f} 10:22 (E13, E14),
\texttt{ec8e206} 11:39 (E15 to E19), \texttt{e38f21e} 11:54 (E20, E21).
No entry can end after the commit that contains it.

\textbf{How it was found.} By me, on 22 November 2026 at 15:35 PST, when
I read the clock of the workspace before writing new entries and
compared the commit times with the log.

\textbf{Provenance.} Each time was written with its entry, in the commit
named above.

\textbf{Why I made it.} I wrote the times without reading a clock; my
estimate of the elapsed time ran ahead of it.

\textbf{The correction.} The intervals of the corrected entries are now
bounded by the times of the commits: an interval ends at the commit that
contains the entry, and a start that I cannot establish is replaced by
the time of the previous commit (entries E13 to E21). The question of C4
was sent after the commit \texttt{e38f21e}. From now on every time in
this record is read from the clock of the workspace when the entry is
written. The same error in two entries of the task record of the folder
of the revision for the Higgs field is corrected there.

\hypertarget{the-decisions-of-interpretation}{%
\section{5. The decisions of
interpretation}\label{the-decisions-of-interpretation}}

\textbf{D1. The folder.} The work is done in the new folder
\texttt{Revision\_Lord\_Gravity}, built on the earlier hypotheses. The
folders of the earlier hypotheses and of the revision for the Higgs
field of the Standard Model stay unchanged. \emph{Reason:} the author
chose the path of the brochon pair bound by the zeugon (Section 2.2),
and the zeugon exists only under the earlier hypotheses; under the
revised hypotheses the massless gauge boson is the photon.

\textbf{D2. The name of the spinor.} The spinor field \(\Psi_{\rm LL}\)
of the earlier hypotheses is called the brochon. Its symbol stays
\(\Psi_{\rm LL}\), because it is the same object (Rule 8).

\textbf{D3. The symbols.} Symbols taken from the earlier work on the
super-loop keep their suffix LL, symbols of the earlier hypotheses their
suffix SF, and every symbol that this folder introduces carries the
suffix LG (Lord gravity), which no other folder of the project uses.
Symbols of Lord's paper are written with the suffix Lord where they are
quoted as his, so that his \(\Gamma_a\), \(S\), \(g\) and the project's
objects of the same letters are never confused.

\textbf{D4. The reading of the paper.} All twenty pages of the paper are
read, and each of its sections is restated in the manuscript in my
words, with its equations re-derived where the manuscript uses them.
Where an equation of the paper, as printed, does not follow from its
definitions, the manuscript says so, shows the corrected form, and a
program verifies both; such findings are recorded in the manuscript as
remarks, not silently repaired.

\textbf{D5. Lord's symbols and his trace.} Where an object of Lord's
paper is the same as an object of the project (his \(\sigma^\mu\),
\(\epsilon\), \(\gamma^\mu\), \(C\) and \(\eta\) are the project's
\(\sigma^{\rm LL}_\mu\), \(\mathsf E_{\rm LL}\),
\(\gamma^\mu_{\rm LL}\), \(C_{\rm LL}\), \(\eta^{\rm LL}\); his
\(\gamma^5\) is \(-\gamma^{\rm ch}_{\rm LL}\)), the manuscript writes
the project's symbol, and a program proves the identity; his other
objects carry the suffix Lord. His ``tr'' in his Section 4 is read as
the trace normalized to 1 on the unit matrix, because only this reading
makes his equations (4.15), (4.16), (4.31) and (4.32) consistent with
each other. \emph{Reason:} Rule 8, and the checks of the manuscript
(Part I, Section 5.3).

\textbf{D6. The conventions of Part III.} Part III works in Lord's
special gauge, in which the Dirac matrices of the frame are the constant
matrices \(\gamma^a_{\rm LL}\) and the tetrad carries the gravitational
field (the vierbein form); it keeps the signature \((+,-,-,-)\) of the
project; it uses the curvature conventions of Landau and Lifshitz, with
which Einstein's equations read
\(G^{\rm LG}_{\mu\nu}=8\pi G_{\rm N}T^{\rm LG}_{\mu\nu}\) with a
positive energy density; and it defines the spin connection
\(\omega^{\rm LG}_{\mu\,ab}\) of the tetrad by the tetrad postulate, so
that the spinor connection
\(\tfrac14\omega^{\rm LG}_{\mu\,ab}\gamma^{ab}_{\rm LL}\) is the spin
connection \(\Gamma^{\rm LL}_\mu\) of GEE with its sign in the earlier
hypotheses and minus Lord's auxiliary field. \emph{Reason:} the
project's metric and the sign of the GEE derivative are fixed by the
earlier hypotheses; Lord's own signs need the corrections of Part I; and
the readers know these conventions from the textbooks. Every sign is
decided by a check (Part III).

\textbf{D7. The source of gravitation.} The energy-momentum tensor of
the brochon is defined by the variation of the action of the earlier
hypotheses with respect to the tetrad. \emph{Reason:} Lord identifies
the source of his Einstein equation with ``a multiple of the
stress-energy tensor'' of matter without defining it, and he abandons
the Lagrangian; the earlier hypotheses have a Lagrangian, and the
variation of the tetrad is the definition that gives a symmetric tensor
whose conservation the field equations guarantee (checked in Part III).

\textbf{D8. The level of description.} Part III states the field
equations of the field and of the mean field, and Einstein's equations
with the expectation value of the energy-momentum tensor as the source;
the brochon in Parts IV and V is a test particle in a given geometry.
\emph{Reason:} the normal order and the momentum cutoff of the earlier
hypotheses have no unique generalization to a curved space-time without
a time-like symmetry; items 5 and 6 of the task need the one-particle
equations, which need neither. The manuscript explains this in a remark
(Section 15.4).

\textbf{D9. Which brochon is put into the gravitational fields.} The
zeugon exists only in the condensed vacuum of the earlier hypotheses,
and there the excitations of the brochon are quasi-particles; by Theorem
L2 of Part III the quasi-particle is, in every gravitational field, a
particle of spin \(\tfrac12\) and mass \(E^{\rm CJT,gap}_{\rm SF}\) (the
gap). Parts IV to VI therefore solve the Dirac equation with this mass,
and give the results for the mass \(m_{\rm LL}\) of the brochon of the
normal vacuum alongside, since the solutions in the Schwarzschild field
depend on the mass only through its product with the mass of the black
hole. \emph{Reason:} the author's path (C2) needs the zeugon. If the
author prefers the brochon of the normal vacuum, the same solutions
serve. \emph{Added in E27:} Theorem L2 gives only one direction (every
solution of the equation of the quasi-particle solves the equation of
the second order of the Dirac particle); the direction that the use of
the Dirac solutions needs is Theorem L3 of Section 18.1 of the
manuscript, added on 22 November 2026.

\textbf{D10. The space-time of a black hole whose mass changes (stage
6).} Stage 6 uses the \emph{ingoing Vaidya space-time},
\(ds^2=f_{\rm LG}\,dv^2-2\,dv\,dr-r^2(d\theta^2+\sin^2\theta\,d\varphi^2)\)
with \(f_{\rm LG}=1-2\mathcal M_{\rm LG}(v)/r\), where the
\emph{advanced time} \(v\) is constant along every light ray that falls
radially inward. It is the exact solution of Einstein's equations whose
source is a radial inflow of \emph{null dust}, radiation without
pressure that moves at the speed of light, with the energy density
\((d\mathcal M_{\rm LG}/dv)/(4\pi r^2)\) (in units \(G_{\rm N}=c=1\)):
for \(d\mathcal M_{\rm LG}/dv>0\) a black hole that grows by absorbing
the inflow. A shrinking black hole is described by the same metric with
\(d\mathcal M_{\rm LG}/dv<0\), an inflow of negative energy: this is the
model of an evaporating black hole near its horizon that W. A. Hiscock
constructed (``Models of evaporating black holes. I'', Phys. Rev.~D 23,
2813 (1981)), in which the negative energy that falls in represents the
reaction of the black hole to the emission of Hawking radiation; like
evaporation itself, it violates the condition that the energy density is
positive. The \emph{apparent horizon} is the sphere
\(r=2\mathcal M_{\rm LG}(v)\) on which the outgoing radial light rays
stop moving outward (inside it they move inward); for a constant mass it
is the horizon. The coframe is \(e^0=\tfrac12(1+f_{\rm LG})\,dv-dr\),
\(e^1=\tfrac12(1-f_{\rm LG})\,dv+dr\), \(e^2=r\,d\theta\),
\(e^3=r\sin\theta\,d\varphi\), regular at every \(r>0\). \emph{Reason:}
the plan (Section 6) asks for a mass that grows by absorbing null dust
and for one that shrinks by radiating it, across the apparent horizon.
The \emph{outgoing} Vaidya space-time, whose mass decreases with the
retarded time because it emits outgoing null dust, describes the
radiation far from the black hole, but its coordinates are singular at
the horizon, so that the field could not be followed across it. The
radiation is therefore represented by its effect at the horizon, the
inflow of negative energy. If the author wants the outgoing metric, it
can be added for the region outside the horizon, joined to the ingoing
one as in Hiscock's model.

\textbf{D11. The brochon pair of stage 7.} The pair that the author
named (C2) is computed as two quasi-particles of the condensed vacuum
with the charges \(+\tfrac12\) and \(-\tfrac12\) of
\(Q_{\rm LL}=T^3_{\rm LL}\), bound by the exchange of the zeugon. By D9
and Theorem L2 of Part III each is a particle of spin \(\tfrac12\) and
mass \(E^{\rm CJT,gap}_{\rm SF}\) whose equation of the second order is
that of a Dirac particle, so that its magnetic moment has the
gyromagnetic factor 2; the results are also given for the mass
\(m_{\rm LL}\) of the brochon of the normal vacuum. The exchange of the
massless zeugon between the charges \(\pm\tfrac12g^{\rm v}_{\rm LL}\)
gives the attractive potential \(-\alpha^{\rm Z}_{\rm LG}/r\) with
\(\alpha^{\rm Z}_{\rm LG}=(g^{\rm v}_{\rm LL})^2/(16\pi)\) (units
\(\hbar=c=1\), the field normalized as usual), at the coupling
\(g^{\rm v}_{\rm LL}=0.7\) of the earlier hypotheses. \emph{Reason:}
this is the pair of C2 in the vacuum where the zeugon exists, with the
masses and charges that Part III derived. Three questions are left open
and stated with every result: (i) whether two quasi-particles of
opposite charge can annihilate into zeugons, since a quasi-particle
mixes particles and antiparticles of the brochon (the levels are given
without and with the term of virtual annihilation, the decay rates only
for the case that annihilation exists); (ii) the change of the coupling
at long distances by the polarization of the condensed vacuum, which
acts like the dielectric constant of a medium and is not computed; (iii)
the value of the cutoff \(\Lambda_{\rm LL}\) in physical units, on which
every prediction in kilograms or electronvolts depends: the earlier
hypotheses fix only ratios to it. On 22 November 2026, about 12:00 PST,
the author was asked which gravitational predictions he means (item 6 of
the task) and whether \(\Lambda_{\rm LL}\) is fixed.

\textbf{D12. The route of stage 7 (my decision, 22 November 2026, after
the author's answer C5).} (1) Every prediction in physical units is
given as a function of \(\Lambda_{\rm LL}\), with a scan from 1 eV to 1
TeV, since the earlier hypotheses fix only ratios to it. (2) The effects
of black holes come in this order: first black holes that carry zeugon
charge (a black hole that absorbs more quasi-particles of one charge
than of the other becomes charged under the zeugon; since the zeugon is
massless, the charge changes the geometry around the black hole; the
charge is limited by the creation of pairs of quasi-particles in the
strong field near the horizon; and it changes the ringing of a black
hole formed in a merger and the size of its shadow); then the tidal
break-up of the pair, refined by the rate of tunnelling along a real
fall. (3) The black holes stay non-rotating: rotation would be needed
only for the growth of clouds of bosons by superradiance, a route not
taken. (4) The pair stays that of D11. (5) Whether two quasi-particles
of opposite charge can annihilate into zeugons is computed from the
mixing of particles and antiparticles in the quasi-particles.
\emph{Reason:} the charged black holes involve both the zeugon and the
spinor and change quantities that are measured; the tidal break-up
continues E20; the other routes need either a rotating black hole or the
physical scale.

\hypertarget{the-plan}{%
\section{6. The plan}\label{the-plan}}

The task is carried out in stages; each is pushed when it is complete.

\begin{enumerate}
\def\labelenumi{\arabic{enumi}.}
\tightlist
\item
  \textbf{Stage 1:} the folder, its document pipeline, this record and
  the README.
\item
  \textbf{Stage 2:} Lord's paper read and understood: Part I of the
  manuscript, with programs that check his conventions and equations
  (Rust, Python, Wolfram).
\item
  \textbf{Stage 3:} the association of the project's \(SL(2,\mathbb C)\)
  with Lord's gauging of the representation of the Clifford algebra:
  Part II, with checks.
\item
  \textbf{Stage 4:} the gravitational coupling of the brochon by Lord's
  covariant derivative: the field equations of the brochon in a curved
  space-time, its energy-momentum tensor, and the gravitational field
  equations with the brochon and the zeugon as sources: Part III, with
  checks.
\item
  \textbf{Stage 5:} the brochon in the Schwarzschild field: separation
  of the field equations, the radial equations, the behaviour on both
  sides of the horizon in coordinates regular there, the quasi-bound
  states and the scattering, integrated with SUNDIALS: Part IV.
\item
  \textbf{Stage 6:} the brochon in the Vaidya metric, for a mass that
  grows by absorbing null dust and for one that shrinks by radiating it,
  across the apparent horizon: Part V.
\item
  \textbf{Stage 7:} the brochon pair bound by the zeugon, and the
  predictions involving the zeugon and the brochon that can be detected
  gravitationally: Part VI.
\end{enumerate}

Added on 22 November 2026 (entry E5), because the rules of the project
require them of every folder:

\begin{itemize}
\tightlist
\item
  \textbf{Stage 4b}, after stage 4: the provenance of the remarks
  (\texttt{P4}), which lists every qualifying remark of the manuscript
  with what it means, why it holds and why it matters, and the registry
  of the symbols (\texttt{P1}), which relates the symbols of this folder
  to those of the other folders; both for Parts I to III, and extended
  with every later part.
\item
  \textbf{Stage 8}, after stage 7: the notebook of the folder (Wolfram
  Language), with every formula and check of the manuscript, its
  remarks, and its provenance (\texttt{P3}).
\end{itemize}

Changed on 22 November 2026 (entry C3): the author instructed to
continue the calculations and the analysis and to stop writing
manuscripts. Stages 5 to 7 are therefore carried out as computations,
each with its programs, data, reports of checks and entries in Section
7; Parts IV to VI of the manuscript, stage 4b and stage 8 wait until the
author asks for documents again, and the rewriting of Parts I to III
(Section 4.1) comes first then. Stage 5 was computed in entries E7 to
E14, stage 6 in entries E15 to E19 (with decision D10 on the space-time
of a shrinking black hole), and the first part of stage 7, the levels of
the pair and its tidal dissociation, in entries E20 and E21 (decision
D11). On 22 November 2026 at 13:51 PST the author asked that the
manuscripts be written again (C5): the rewriting of Parts I to III
(Section 4.1) comes first, then Parts IV to VI from stages 5 to 7, in
the friendlier style that he asked for; the rest of stage 7 follows
decision D12. Parts I to III were rewritten in entries E22 to E24, and
Part IV was written from stage 5 in entries E26 and E27. Next: Part V
from stage 6; then the rest of stage 7 (D12) and Part VI.

The plan is changed only with a record in Section 7.

\hypertarget{log-of-the-requests-the-work-the-errors-and-the-changes}{%
\section{7. Log of the requests, the work, the errors and the
changes}\label{log-of-the-requests-the-work-the-errors-and-the-changes}}

Entries C are the author's messages that concern this folder; entries E
are the work done. The verdicts and numbers that the entries quote are
written by the programs into the fragment files, and they are those of
the current version of each program: a program that gained checks after
an entry (for example \texttt{lg\_schw} in E14 and E27) reports the
later number of checks in the earlier entry as well.

\begin{longtable}[]{@{}
  >{\raggedright\arraybackslash}p{(\columnwidth - 4\tabcolsep) * \real{0.0610}}
  >{\raggedright\arraybackslash}p{(\columnwidth - 4\tabcolsep) * \real{0.1829}}
  >{\raggedright\arraybackslash}p{(\columnwidth - 4\tabcolsep) * \real{0.7561}}@{}}
\toprule\noalign{}
\begin{minipage}[b]{\linewidth}\raggedright
entry
\end{minipage} & \begin{minipage}[b]{\linewidth}\raggedright
time (PST)
\end{minipage} & \begin{minipage}[b]{\linewidth}\raggedright
what
\end{minipage} \\
\midrule\noalign{}
\endhead
\bottomrule\noalign{}
\endlastfoot
C1 & 21 Nov 2026, 18:02 & the task C23 (Section 2.2), to follow the
stages of the revision of the hypotheses \\
C2 & 21 Nov 2026, 18:45 & the path of the task: the brochon pair bound
by the exchange of the zeugon; binding by photons ``extremely unlikely''
and ``not worth studying''; the name brochon \\
E1 & 22 Nov 2026, 04:35 to 04:40 & stage 1: the folder
\texttt{Revision\_Lord\_Gravity} with its document pipeline
(\texttt{code/scripts/}: \texttt{assemble\_docs.py},
\texttt{build\_docs.sh}, \texttt{doc\_header.tex},
\texttt{keep\_tables.py}, \texttt{check\_layout.py},
\texttt{check\_pages.py}, \texttt{check\_self\_contained.py}, adapted
from the folder of the revision of the hypotheses with the folder name
and the running head changed), this record, the README, and the checksum
of Lord's paper in \texttt{provenance/inputs\_sha256.txt}. The paper was
read completely (all twenty pages, the journal pages 423 to 442) \\
E2 & 22 Nov 2026, 04:45 to 05:02 & stage 2: Part I of the manuscript
\texttt{lord\_gravity.md} (Lord's five sections restated with the
equations that the task uses; Section 7, the equations whose printed
form does not follow from his definitions: the sign of the inhomogeneous
term of (4.14), (4.31), (4.16), (4.32), the factor 2 of the connection
(4.20), (5.2), the coefficient of the curvature (4.47), (4.48), and the
reading of ``any linearly independent set'' of Hermitian matrices). The
Rust workspace \texttt{rust/} (library \texttt{lgcore}, with modules
copied from the earlier folder and the new module \texttt{lord.rs};
program \texttt{lg\_lord}), the Python programs
\texttt{python/lgcommon.py} and \texttt{python/lg\_lord.py}, the Wolfram
script \texttt{wolfram/source/lg\_lord.wls} (run on the author's
computer by the watcher, request 325), and
\texttt{code/scripts/write\_lg\_lord\_fragments.py}. Results: Rust ALL
PASS (46 of 46), Python ALL PASS (17 of 17), Wolfram ALL PASS (7 of 7).
The same programs already check the dictionary of Part II (checks P-1 to
P-8) \\
E3 & 22 Nov 2026, 05:03 to 05:08 & stage 3: Part II of
\texttt{lord\_gravity.md} (Sections 9 to 12): the project's objects
(bi-quaternions, M31, frames, the two actions of a unit bi-quaternion,
the brochon, GEE and the zeugon), the dictionary between Lord's objects
and the project's (Table 10.1, every line with its check), Theorem L1
(the continuous changes of representation that keep the Weyl form are
exactly the unit bi-quaternions; new checks P-9 in Rust and YL-9 in
Python), and the association: Lord's representations of the Clifford
algebra are the orthonormal frames of M31, his change of representation
is the action of a unit bi-quaternion of the super-loop on frames, his
gauge group is the \(SL(2,\mathbb C)\) factor of GEE, and his auxiliary
fields are minus its spin connection; what the gauging adds is the
dynamics of that connection. The provenance of the manuscript became
Section 13 and its symbols Section 14 \\
E4 & 22 Nov 2026, 05:08 to 06:25 & stage 4: Part III of
\texttt{lord\_gravity.md} (Sections 13 to 17): the tetrad as Lord's
reference vectors (four Hermitian bi-quaternions at every point whose
Minkowski products are the inverse metric); Lord's change of
representation and the rotation of the tetrad as one local
\(SL(2,\mathbb C)\); the Christoffel symbols and the curvature in the
conventions of D6; the spin connection by the tetrad postulate and by
the anholonomy, the spinor connection
\(\Gamma^{\rm LL}_\mu=\tfrac14\omega^{\rm LG}_{\mu\,ab}\gamma^{ab}_{\rm LL}\),
its curvature \(\tfrac14R^{\rm LG}_{ab\mu\nu}\gamma^{ab}_{\rm LL}\), and
Lord's objects realized by the tetrad; the spinor connection of the
Schwarzschild space-time in the static and in the freely falling
(Painlevé-Gullstrand) tetrad, and the Vaidya space-times; the brochon's
covariant derivative, Lagrangian, field equation, currents, the square
of its Dirac operator with \(-\tfrac14R^{\rm LG}\), the conformal
invariance of the massless brochon and Volkov's solution in a wave of
the zeugon; Theorem L2 (the quasi-particle of the superfluid is a
particle of spin \(\tfrac12\) and mass \(E^{\rm CJT,gap}_{\rm SF}\) in
every gravitational field); the energy-momentum tensor from the
variation of the tetrad, the tensors of the gauge fields and of the
zeugon, and Einstein's equations (D7, D8, D9). \\
E5 & 22 Nov 2026, 05:08 to 06:25 & the programs and the records of stage
4 (E4). Programs: the module \texttt{rust/lgcore/src/curved.rs} (added;
the other modules unchanged), the program \texttt{lg\_curved} (crate
\texttt{rust/lgcurved}), \texttt{python/lg\_curved.py},
\texttt{wolfram/source/lg\_curved.wls} and
\texttt{code/scripts/write\_lg\_curved\_fragments.py}. Results: Rust ALL
PASS (57 of 57), Python ALL PASS (19 of 19), Wolfram ALL PASS (6 of 6).
The Wolfram script reached its form in four runs on the author's
computer (watcher requests 326 to 329): the first two failed the checks
WG-1 to WG-3 because the iterator \texttt{r} of the sums of the
Christoffel and Riemann tensors replaced the coordinate \texttt{r}
inside the sums (the iterator is now \texttt{rr}); the first also spent
18 minutes in \texttt{Simplify} on the resulting meaningless
expressions, and from the second on the identities of the conformally
flat metric and of the variation of the tetrad are decided by expansion
to the canonical polynomial form, and the statements that something is
not zero by its value for explicit functions at a rational point; the
third failed only the curvature of the static Schwarzschild tetrad,
whose square roots \texttt{Simplify} did not remove, and the fourth
checks that curvature for every static, spherically symmetric tetrad
\(\mathrm{diag}(h_{\rm LG}(r),1/h_{\rm LG}(r),r,r\sin\theta)\) with a
symbolic function \(h_{\rm LG}\), which contains the Schwarzschild
tetrad and has no square roots (6 of 6, in 0.7 s; the version in the
repository). Equation (11.1) was rewritten with the spin connection
\(\omega^{\rm LG}\) of Part III; the former \(\omega\) was Lord's
connection of the vectors \(W\), its negative, and both forms are now
stated (Section 4). The provenance of the manuscript became Section 18
and its symbols Section 19. The plan was extended by stages 4b and 8
(Section 6) \\
C3 & 22 Nov 2026, 07:26 & the author's instruction, given about the
article of another folder of the project (the revision for the Higgs
field of the Standard Model) and extended by him to all of my
manuscripts: a document must not send its reader to earlier work that
almost nobody can read, but must contain all the material that it needs;
it must give concrete and explicit examples of the objects that it names
and define every term that it uses, for a student who wants to
understand and to learn rather than to repeat jargon; the calculations
and the analysis are to be continued, but no more manuscripts are to be
written; and a handoff file for another assistant is to be written, in
case he must end this work. What this means for this folder is recorded
in Section 4.1 and in the plan (Section 6) \\
E6 & 22 Nov 2026, 09:40 to 09:52 & the audit of this folder after C3
(Section 4.1) and its corrections: the sentence of Section 2.3; the
check \texttt{code/scripts/check\_self\_contained.py}, whose report of
HCK-H1 now states what it searched, and whose new check HCK-H2 fails,
listing the 30 lines of the source of the manuscript, and the 30 of the
Markdown built from it, that contain the words ``earlier hypotheses'' or
``earlier work''; and the plan (Section 6). No manuscript text \\
E7 & 22 Nov 2026, 08:52 to 09:40 & stage 5, the brochon in the
Schwarzschild field, as computations with checks; no manuscript text
(C3). \emph{The setting.} A Dirac particle of mass \(m_{\rm LL}\), one
isospin component of the brochon (by D9 the same solutions describe the
quasi-particle of the superfluid, with \(m_{\rm LL}\) replaced by its
gap \(E^{\rm CJT,gap}_{\rm SF}\)), with no field of the zeugon, in the
Schwarzschild space-time of a mass \(M_{\rm LG}\): the space-time
outside a spherical mass that has no charge and does not rotate,
continued inward through its horizon, the sphere
\(r=2\mathcal M_{\rm LG}\) from inside which nothing can come out. Units
\(G_{\rm N}=c=\hbar=1\) and
\(\mathcal M_{\rm LG}=G_{\rm N}M_{\rm LG}=1\), so that the horizon is at
\(r=2\) and the solutions depend on the two masses only through the
number \(\mathcal M_{\rm LG}m_{\rm LL}\), which is
\(G_{\rm N}M_{\rm LG}m_{\rm LL}/(\hbar c)\) in ordinary units and plays
the part that the fine-structure constant plays in the hydrogen atom.
Two tetrads: the static one, with the coframe \(f_{\rm LG}^{1/2}dt\),
\(f_{\rm LG}^{-1/2}dr\), \(r\,d\theta\), \(r\sin\theta\,d\varphi\),
\(f_{\rm LG}=1-2/r\), which is singular at the horizon; and the freely
falling one of Painlevé and Gullstrand, with the coframe
\(d\tau_{\rm LG}\), \(dr+w_{\rm LG}\,d\tau_{\rm LG}\), \(r\,d\theta\),
\(r\sin\theta\,d\varphi\), where \(w_{\rm LG}=(2/r)^{1/2}\) is the speed
of an observer who falls from rest at infinity and \(\tau_{\rm LG}\) is
his proper time; it is regular at the horizon and serves on both sides
of it. \emph{The separation.} For a mode of frequency
\(\omega^{\rm S}_{\rm LG}\) (the field depends on the time as
\(e^{-i\omega^{\rm S}_{\rm LG}\tau_{\rm LG}}\)), total angular momentum
\(j_{\rm LG}=\tfrac12,\tfrac32,\dots\) and component \(m_{j,{\rm LG}}\)
of it along the axis, the Dirac equation of Part III separates: the
spinor is
\(r^{-1}(\sin\theta)^{-1/2}e^{-i\omega^{\rm S}_{\rm LG}\tau_{\rm LG}+im_{j,{\rm LG}}\varphi}\)
times a constant \(4\times4\) matrix acting on four products of one of
two angular functions \(\alpha_{\rm LG}(\theta)\),
\(\beta_{\rm LG}(\theta)\) with one of two radial functions
\(\mathcal P_{\rm LG}(r)\), \(\mathcal Q_{\rm LG}(r)\). The angular
functions solve
\(\alpha_{\rm LG}'=(m_{j,{\rm LG}}/\sin\theta)\,\alpha_{\rm LG}-\lambda_{\rm LG}\beta_{\rm LG}\)
and
\(\beta_{\rm LG}'=-(m_{j,{\rm LG}}/\sin\theta)\,\beta_{\rm LG}+\lambda_{\rm LG}\alpha_{\rm LG}\),
which have solutions regular at both poles for
\(\lambda_{\rm LG}=\pm(j_{\rm LG}+\tfrac12)\); for example, for
\(j_{\rm LG}=m_{j,{\rm LG}}=\tfrac12\),
\(\alpha_{\rm LG}=\sin^{1/2}(\theta/2)\cos^{3/2}(\theta/2)\) and
\(\beta_{\rm LG}=\lambda_{\rm LG}\sin^{3/2}(\theta/2)\cos^{1/2}(\theta/2)\).
The radial functions solve two coupled linear equations of the first
order (in the static tetrad, with \(f_{\rm LG}^{-1/4}\) in the
prefactor, two equations equivalent to them). The constant
\(\lambda_{\rm LG}\) is minus the number \(\kappa^{\rm H}_{\rm LG}\) of
Dirac's theory of the hydrogen atom that tells apart the two states of
the same \(j_{\rm LG}\): \(\kappa^{\rm H}_{\rm LG}=-1\) for the states
\(S_{1/2}\), \(+1\) for \(P_{1/2}\), \(-2\) for \(P_{3/2}\) and \(+2\)
for \(D_{3/2}\) (found from the levels of small
\(\mathcal M_{\rm LG}m_{\rm LL}\): \(\lambda_{\rm LG}=1\) has a level
with \(n_{\rm LG}=1\), \(\lambda_{\rm LG}=-1\) has none). Here
\(n_{\rm LG}\,\ell_{j}\) is the notation of the hydrogen atom: the
principal quantum number \(n_{\rm LG}\), the orbital angular momentum
\(\ell_{\rm LG}\) written S, P, D for 0, 1, 2, and \(j_{\rm LG}\) as the
index \\
E8 & 22 Nov 2026, 08:52 to 09:40 & stage 5, the horizon and the interior
(E7). The radial equations of the freely falling tetrad have a regular
singular point at the horizon: a point at which the coefficients of the
equations have a simple pole, so that the solutions near it behave as
powers \((r-2)^{s_{\rm LG}}\), where the possible values of
\(s_{\rm LG}\), the exponents, are the eigenvalues of the residue of the
matrix of coefficients. The exponents are \(s_{\rm LG}=0\), a solution
that is analytic at the horizon (it has a convergent power series in
\(r-2\), whose radius of convergence 2 is the distance to \(r=0\)): the
wave that falls in, which continues through the horizon unchanged; and
\(s_{\rm LG}=4i\mathcal M_{\rm LG}\omega^{\rm S}_{\rm LG}-\tfrac12\):
the wave that comes out of the horizon, singular there, which no
physical solution contains. Inside the horizon the radius decreases
along every path of the particle; near \(r=0\), where the curvature
diverges, both radial functions of the solution that fell in behave as
\(r^{1/4}\) (the residue there is \(\tfrac14\) times the unit matrix),
so that the spinor grows as \(r^{-3/4}\). For a real frequency the
radial current
\(J_{\rm LG}=(1+w_{\rm LG})\lvert\mathcal Q_{\rm LG}\rvert^2-(1-w_{\rm LG})\lvert\mathcal P_{\rm LG}\rvert^2\)
(the flux of the conserved current of the brochon into the sphere of
radius \(r\), divided by the integral of
\(\lvert\alpha_{\rm LG}\rvert^2+\lvert\beta_{\rm LG}\rvert^2\) over the
sphere) is the same at every radius, inside the horizon as outside:
exactly, by computer algebra, and to \(1.0\times10^{-9}\) (relative)
along the computed solutions from \(r=2500\) down to
\(r=1.5\times10^{-12}\) \\
E9 & 22 Nov 2026, 08:52 to 09:40 & stage 5, the quasi-bound states (E7).
A \emph{quasi-bound state} is a solution that falls into the horizon and
decays at infinity; it exists only at complex frequencies,
\(\omega^{\rm S}_{\rm LG}\) with \({\rm Re}\,\omega^{\rm S}_{\rm LG}\)
the energy of the level and \({\rm Im}\,\omega^{\rm S}_{\rm LG}<0\): the
amplitude decays as
\(e^{-\lvert{\rm Im}\,\omega^{\rm S}_{\rm LG}\rvert\tau_{\rm LG}}\)
because the black hole absorbs the particle. For small
\(\mathcal M_{\rm LG}m_{\rm LL}\) the levels approach those of a
hydrogen atom,
\(m_{\rm LL}\big(1-(\mathcal M_{\rm LG}m_{\rm LL})^2/(2n_{\rm LG}^2)\big)\);
at \(\mathcal M_{\rm LG}m_{\rm LL}=0.05\) every computed level lies
within \(2.75\) \((\mathcal M_{\rm LG}m_{\rm LL})^4\,m_{\rm LL}\) of it.
Six states were followed from \(\mathcal M_{\rm LG}m_{\rm LL}=0.05\) to
\(0.4\) (data in \texttt{data/lg\_schw\_qbs.csv}). The values of
\(\omega^{\rm S}_{\rm LG}/m_{\rm LL}\) at
\(\mathcal M_{\rm LG}m_{\rm LL}=0.1\): \(1S_{1/2}\):
\(0.9946882775-2.79\times10^{-5}\,i\); \(2S_{1/2}\):
\(0.9986985113-3.50\times10^{-6}\,i\); \(2P_{1/2}\):
\(0.9987296936-5.80\times10^{-9}\,i\); \(2P_{3/2}\):
\(0.9987396568-6.03\times10^{-12}\,i\); \(3P_{3/2}\):
\(0.9994402017-2.13\times10^{-12}\,i\); \(3D_{3/2}\):
\(0.9994417316\ (\lvert{\rm Im}\rvert<10^{-12})\); at \(0.3\):
\(1S_{1/2}\): \(0.9454639375-0.0296\,i\); \(2S_{1/2}\):
\(0.9861973324-4.54\times10^{-3}\,i\); \(2P_{1/2}\):
\(0.9861834925-4.08\times10^{-4}\,i\); \(2P_{3/2}\):
\(0.9877468037-3.30\times10^{-6}\,i\); \(3P_{3/2}\):
\(0.9945885005-1.30\times10^{-6}\,i\); \(3D_{3/2}\):
\(0.9947595601-3.29\times10^{-9}\,i\); the state \(1S_{1/2}\) at
\(0.4\): \(0.930061-0.0701\,i\), far from its hydrogen-like value
\(0.9200\). The amplitude of the state \(1S_{1/2}\) at
\(\mathcal M_{\rm LG}m_{\rm LL}=0.1\) decays by the factor \(e\) in the
time \(3.59\times10^{5}\) \(G_{\rm N}M_{\rm LG}/c^3\). The decay rate
\(\lvert{\rm Im}\,\omega^{\rm S}_{\rm LG}\rvert\) falls with
\(\mathcal M_{\rm LG}m_{\rm LL}\) as a power that grows with
\(\ell_{\rm LG}\) and \(j_{\rm LG}\). Between neighbouring computed
masses (the smallest at which the roots resolve
\(\lvert{\rm Im}\,\omega^{\rm S}_{\rm LG}\rvert\ge10^{-12}m_{\rm LL}\))
the local powers are \(1S_{1/2}\): 6.03
(\(\mathcal M_{\rm LG}m_{\rm LL}\) from 0.05 to 0.1); \(2S_{1/2}\): 6.03
(\(\mathcal M_{\rm LG}m_{\rm LL}\) from 0.05 to 0.1); \(2P_{1/2}\): 7.92
(\(\mathcal M_{\rm LG}m_{\rm LL}\) from 0.05 to 0.1); \(2P_{3/2}\):
10.87 (\(\mathcal M_{\rm LG}m_{\rm LL}\) from 0.1 to 0.15);
\(3P_{3/2}\): 10.90 (\(\mathcal M_{\rm LG}m_{\rm LL}\) from 0.1 to
0.15); \(3D_{3/2}\): 14.79 (\(\mathcal M_{\rm LG}m_{\rm LL}\) from 0.2
to 0.25); none of them is the limiting power of small masses, because at
these masses the corrections of relative order
\((\mathcal M_{\rm LG}m_{\rm LL})^2\) are still large. The limiting
powers are 5, 5, 7, 9, 9 and 11 (entry E13; until 22 November 2026 this
entry said otherwise, Section 4.2) \\
E10 & 22 Nov 2026, 08:52 to 09:40 & stage 5, the scattering (E7). For a
real frequency above the mass, a partial wave (a wave of definite
\(j_{\rm LG}\), \(m_{j,{\rm LG}}\) and \(\lambda_{\rm LG}\)) that comes
in from infinity is partly absorbed by the black hole and partly
scattered back; the absorbed fraction of its flux is the \emph{greybody
factor} \(\mathcal A_{\rm LG}\) of the partial wave (the name says that
the black hole absorbs like a body that is not perfectly black). It is
computed from the solution that falls in, decomposed at a large radius
into the waves that go out and come in (data in
\texttt{data/lg\_schw\_grey.csv}). For the massless brochon at
\(\mathcal M_{\rm LG}\omega^{\rm S}_{\rm LG}=0.25\), \(0.4\) and
\(0.6\): \(\mathcal A_{\rm LG}=\) \(0.873387\), \(0.998219\),
\(0.999993\) for \(\lambda_{\rm LG}=\pm1\), the two signs agreeing to
\(2.2\times10^{-16}\), as they must, since for \(m_{\rm LL}=0\) a
substitution of the two radial functions (check YS-9) turns the
equations of \(\lambda_{\rm LG}\) into those of \(-\lambda_{\rm LG}\)
and keeps the current; and \(0.009528\), \(0.632114\), \(0.998753\) for
\(\lambda_{\rm LG}=2\): at the lowest of these frequencies almost the
whole wave of the larger angular momentum is scattered back. With the
mass, \(\mathcal M_{\rm LG}m_{\rm LL}=0.2\) at
\(\mathcal M_{\rm LG}\omega^{\rm S}_{\rm LG}=0.25\), the two signs
differ: \(0.9017\) for \(\lambda_{\rm LG}=1\) and \(0.4763\) for
\(\lambda_{\rm LG}=-1\) \\
E11 & 22 Nov 2026, 08:52 to 09:52 & stage 5, the programs (E7 to E10).
Rust: the crate \texttt{rust/lgsundials} (the integrator CVODE of
SUNDIALS 7.8.0 in pure Rust, a copy of the module of the folder of the
revision for the Higgs field) and its C interface
\texttt{rust/lgsundials\_c}, and the program \texttt{lg\_schw} (crate
\texttt{rust/lgschw}; ALL PASS (21 of 21)): the separation checked by
the Dirac operator of Part III at random points, the series at the
horizon, both sides of the horizon, the interior, the quasi-bound states
(the solution that falls in and the solution that decays at infinity,
integrated by CVODE toward each other and matched by the determinant of
the two solution vectors, the root found by a complex secant iteration;
the outer radius 1.3 times larger changes no frequency by more than
\(4.3\times10^{-14}\), relative) and the greybody factors. Python:
\texttt{python/lg\_schw.py} (ALL PASS (14 of 14)), with the module
\texttt{python/lgsundials.py}: the separation, the relation of the two
tetrads, the exponents and the current by computer algebra (sympy), and
the quasi-bound states and greybody factors recomputed from its own
start (a Taylor polynomial at \(r=2.002\)), with CVODE and Muller's
method; the frequencies agree with the Rust values to
\(1.7\times10^{-14}\) \(m_{\rm LL}\) and the greybody factors to
\(1.7\times10^{-9}\). Wolfram: \texttt{wolfram/source/lg\_schw.wls} (ALL
PASS (9 of 9), Wolfram 15.0.1 on the author's computer through the
watcher, requests 337, 338 and 341; WS-9 was added with E13): the
separation exactly, with the spin connection computed from each tetrad,
the relation, the exponents, the current and the angular functions of
\(j_{\rm LG}=\tfrac12\) exactly, and three quasi-bound states and three
greybody factors in arithmetic of 32 digits with the integrator NDSolve
of the Wolfram Language (SUNDIALS offers no arithmetic of more than
double precision; the computation serves only as a check), which agree
with the Rust values to \(7.3\times10^{-14}\) \(m_{\rm LL}\) and
\(6.6\times10^{-10}\). The writer of the fragments:
\texttt{code/scripts/write\_lg\_schw\_fragments.py} \\
E12 & 22 Nov 2026, 09:00 to 09:52 & stage 5, the corrections before any
result was used (E11). (1) Rust: the first criterion for a quasi-bound
state, the ratio of the growing to the decaying part of the solution
from the horizon at the outer radius, gave spurious roots near
\(\omega^{\rm S}_{\rm LG}=m_{\rm LL}\), where both parts vary slowly; it
was replaced by the matching of E11. (2) The continuation of a state
from one mass to the next jumped to a neighbouring state until the start
was extrapolated (the distance of \({\rm Re}\,\omega^{\rm S}_{\rm LG}\)
from the hydrogen-like level scaled with
\((\mathcal M_{\rm LG}m_{\rm LL})^4\),
\({\rm Im}\,\omega^{\rm S}_{\rm LG}\) with
\((\mathcal M_{\rm LG}m_{\rm LL})^6\)). (3) The check that
\({\rm Im}\,\omega^{\rm S}_{\rm LG}<0\) failed for the state
\(3D_{3/2}\), whose computed imaginary parts at the smallest masses were
of the order \(10^{-16}m_{\rm LL}\) and of either sign, below the
resolution of the computation; the check now requires
\({\rm Im}\,\omega^{\rm S}_{\rm LG}<0\) only where
\(\lvert{\rm Im}\,\omega^{\rm S}_{\rm LG}\rvert\ge10^{-12}m_{\rm LL}\)
and counts the other roots (5 of 48). (4) The checks of the current use
CVODE with the relative tolerance \(10^{-13}\) and accept \(10^{-7}\),
ten times the global error, of the order \(10^{-8}\), that CVODE reaches
over these integrations. (5) Python: the relation of the two tetrads was
first evaluated by the adaptive numerical evaluation of sympy, which
returned \(10^{-180}\) (a zero within its limit of precision); it is now
evaluated in arithmetic of 40 digits (\(5\times10^{-41}\)). The
comparison of the exponent
\(4i\mathcal M_{\rm LG}\omega^{\rm S}_{\rm LG}-\tfrac12\) first failed
because sympy wrote the same number in another form; it now subtracts
the two. The first version integrated with the Runge-Kutta method DOP853
of scipy (a library method, not a hand-written one); to follow the
instruction that every numerical integration uses SUNDIALS it now uses
CVODE through \texttt{rust/lgsundials\_c} (the run with DOP853 gave the
same verdict, with the frequencies agreeing to
\(7.2\times10^{-14}m_{\rm LL}\) instead of \(1.7\times10^{-14}\)
\(m_{\rm LL}\)). The check YS-9 (E10) was added after the first run. (6)
Wolfram: request 337 failed WS-3 and WS-6, because \texttt{Simplify} did
not remove the nested square roots (the residuals of WS-3 were zero but
printed unsimplified; WS-6 printed only False); request 338 decides both
as polynomial identities, in \(a=(1-w_{\rm LG})^{1/2}\) and
\(b=(1+w_{\rm LG})^{1/2}\) with \(a^2+b^2=2\), and in \(\sin(\theta/2)\)
and \(\cos(\theta/2)\) with the sum of squares 1, and passes 8 of 8 \\
E13 & 22 Nov 2026, 09:52 to 10:22 & stage 5, the quasi-bound states at
small masses, computed with checks; no manuscript text (C3). \emph{The
decay rate from the charge.} For every solution of the radial equations
with a complex frequency,
\(dJ_{\rm LG}/dr=2\,{\rm Im}\,\omega^{\rm S}_{\rm LG}\,(\lvert\mathcal P_{\rm LG}\rvert^2+\lvert\mathcal Q_{\rm LG}\rvert^2)\)
exactly (proved by computer algebra, checks YS-10 and WS-9): the charge
of the brochon between the horizon and a sphere changes in time only
through the fluxes through the two. A quasi-bound state decays at
infinity, where \(J_{\rm LG}\) therefore vanishes, and \(J_{\rm LG}=2\)
at the horizon for the normalization \(\mathcal Q_{\rm LG}(2)=1\); hence
\({\rm Im}\,\omega^{\rm S}_{\rm LG}=-1/N_{\rm LG}\), where
\(N_{\rm LG}\) is the integral of
\(\lvert\mathcal P_{\rm LG}\rvert^2+\lvert\mathcal Q_{\rm LG}\rvert^2\)
from \(r=2\) to infinity. At the 30 quasi-bound states of E9 with
\(\lvert{\rm Im}\,\omega^{\rm S}_{\rm LG}\rvert\ge10^{-7}m_{\rm LL}\)
the two sides agree to \(5.5\times10^{-8}\) (relative), at the other 18
to \(1.0\times10^{-14}\) \(m_{\rm LL}\) (LS-10). Evaluated at the real
frequency \({\rm Re}\,\omega^{\rm S}_{\rm LG}\), the same formula gives
the decay rate to the first order in
\({\rm Im}\,\omega^{\rm S}_{\rm LG}\) divided by the distance to the
neighbouring levels, in double precision however small the rate is; at
\(\mathcal M_{\rm LG}m_{\rm LL}=0.05\) it agrees with the resolved roots
to \(1.3\times10^{-5}\) (relative; LS-13). \emph{Small masses.} At
\(\mathcal M_{\rm LG}m_{\rm LL}=0.005\), \(0.01\), \(0.02\), \(0.04\)
and \(0.05\) (data in \texttt{data/lg\_schw\_small.csv}) the six states
were compared with the small-mass formulas that S. R. Dolan and D.
Dempsey give for a Dirac field on a rotating black hole, of which the
Schwarzschild field is the case without rotation (``Bound states of the
Dirac equation on Kerr spacetime'', arXiv:1504.03190, equations (94) to
(96), the last two after Ternov and co-workers). Writing \(n,\ell,j,S\)
for \(n_{\rm LG},\ell_{\rm LG},j_{\rm LG}\) and
\(S_{\rm LG}=j_{\rm LG}-\ell_{\rm LG}\), and
\(y_{\rm LG}=\mathcal M_{\rm LG}m_{\rm LL}\): the energy is
\({\rm Re}\,\omega^{\rm S}_{\rm LG}/m_{\rm LL}=1-y_{\rm LG}^2/(2n^2)+F^{\rm H}_{\rm LG}y_{\rm LG}^4+\dots\)
with
\(F^{\rm H}_{\rm LG}=\big(\tfrac{15}8-\tfrac{3n}{2j+1}-\tfrac{3n}{2\ell+1}\big)/n^4\),
the gravitational counterpart of the fine structure of the hydrogen
atom; and the decay rate is
\(-{\rm Im}\,\omega^{\rm S}_{\rm LG}/m_{\rm LL}=a^{\rm H}_{\rm LG}\,y_{\rm LG}^{4\ell+2S+4}\)
with
\(a^{\rm H}_{\rm LG}=\frac{(n+\ell)!}{n^{4+2\ell}(2\ell)!\,(2\ell+1)!\,(n-\ell-1)!}(2\ell+1)^{-2S}\prod_{p=1}^{j+1/2}\big(1+\tfrac{16y_{\rm LG}^2}{(p-1/2)^2}\big)\).
The computed decay rate is \(1.0027\) times the formula for \(1S_{1/2}\)
and \(1.0010\) times for \(3D_{3/2}\) at \(y_{\rm LG}=0.005\) (\(1.179\)
for \(1S_{1/2}\) at \(y_{\rm LG}=0.05\)); the computed coefficient of
\(y_{\rm LG}^4\) in the energy is \(1.048\) times \(F^{\rm H}_{\rm LG}\)
for \(1S_{1/2}\) and \(1.0025\) times for \(3D_{3/2}\) at
\(y_{\rm LG}=0.05\). Extrapolated to \(y_{\rm LG}=0\), both ratios are 1
for all six states, to \(3.4\times10^{-4}\) for the energy and
\(7.2\times10^{-4}\) for the decay rate (LS-11, LS-12), by two
extrapolations: Richardson's rule for corrections in even powers of
\(y_{\rm LG}\), and a fit with a linear and a quadratic term, since for
the S states the corrections contain odd powers (the differences of the
ratios grow by factors near 3.7, not 4, when \(y_{\rm LG}\) doubles);
the two limits lie on opposite sides of 1. The local powers of the decay
rate between \(y_{\rm LG}=0.005\) and \(0.01\) are \(1S_{1/2}\): 5.017
(limit 5); \(2S_{1/2}\): 5.017 (limit 5); \(2P_{1/2}\): 7.014 (limit 7);
\(2P_{3/2}\): 9.012 (limit 9); \(3P_{3/2}\): 9.012 (limit 9);
\(3D_{3/2}\): 11.012 (limit 11). Two examples of the rates:
\({\rm Im}\,\omega^{\rm S}_{\rm LG}/m_{\rm LL}=\)
\(-3.14\times10^{-12}\) for \(1S_{1/2}\) at \(y_{\rm LG}=0.005\) and
\(-3.21\times10^{-27}\) for \(3D_{3/2}\) at \(y_{\rm LG}=0.01\). The
agreement is an independent confirmation of E9, and, since the formulas
depend on \(\ell\) and \(j\), of the identification
\(\lambda_{\rm LG}=-\kappa^{\rm H}_{\rm LG}\) (E7) \\
E14 & 22 Nov 2026, 09:52 to 10:22 & the programs of E13 and the
corrections on the way. Rust: the program \texttt{lg\_schw} gained the
decay rate from the charge (the function \texttt{decay\_from\_norm}: the
integral of the norm along the CVODE integrations, and over the first
half unit of radius from the horizon series by Gauss-Legendre
quadrature), the small-mass formulas, the checks LS-10 to LS-13 and the
data file \texttt{data/lg\_schw\_small.csv} (ALL PASS (21 of 21)).
Python: \texttt{python/lg\_schw.py} gained YS-10 (the identity, by
sympy), YS-11 (the decay rates of the six states at \(y_{\rm LG}=0.01\)
recomputed with its own start and integrals; agreement
\(9.7\times10^{-11}\)) and YS-12 (the identity at two roots;
\(3.5\times10^{-11}\)) (ALL PASS (14 of 14)). Wolfram:
\texttt{wolfram/source/lg\_schw.wls} gained WS-9 (the identity, exactly;
request 341; ALL PASS (9 of 9)). Corrections: (1) the first plan
computed the quasi-bound states at small masses by the matching of E11
in arithmetic of 50 digits in the Wolfram Language, with a script that
is not kept in the repository; request 339 was rejected by the watcher,
which accepts only paths as arguments, and request 340 failed: the
integrator NDSolve stopped near the horizon (``error test failure'') at
the precision goal of 40 digits in the variable \(\ln r\), and the root
finder then wandered (495 evaluations, no root). The decay rate from the
charge, which needs no more than double precision, replaced it. (2) The
first version of LS-11 and LS-12 extrapolated only by Richardson's rule
for even powers, which left a difference of \(3\times10^{-4}\) from 1
for the S states; since their corrections contain odd powers, the checks
now use both extrapolations \\
E15 & 22 Nov 2026, 10:22 to 11:39 & stage 6, the equations, computed
with checks; no manuscript text (C3). The space-time and the coframe are
those of decision D10; units \(G_{\rm N}=c=\hbar=1\), lengths and times
in units of the mass of the black hole. \emph{The separation.} With the
angular functions \(\alpha_{\rm LG}(\theta)\),
\(\beta_{\rm LG}(\theta)\) of stage 5 and the spinor
\(r^{-1}(\sin\theta)^{-1/2}e^{im_{j,{\rm LG}}\varphi}S_{\rm H}\,(-i\mathcal Q^{\rm V}_{\rm LG}\alpha_{\rm LG},\,i\mathcal P^{\rm V}_{\rm LG}\beta_{\rm LG},\,\mathcal P^{\rm V}_{\rm LG}\alpha_{\rm LG},\,\mathcal Q^{\rm V}_{\rm LG}\beta_{\rm LG})\),
where \(S_{\rm H}=\mathrm{diag}(H_{\rm LG},H_{\rm LG})\) with
\(H_{\rm LG}=i(\sigma_1+\sigma_3)/\sqrt2\) is the constant matrix of
stage 5 (E7; \(\sigma_1\) and \(\sigma_3\) are Pauli matrices:
\(\sigma_1\) has the elements 1 off the diagonal and 0 on it,
\(\sigma_3=\mathrm{diag}(1,-1)\)) and
\(\mathcal P^{\rm V}_{\rm LG}(v,r)\),
\(\mathcal Q^{\rm V}_{\rm LG}(v,r)\) depend on time and radius, the
Dirac equation of Part III with the mass \(m_{\rm LL}\) is equivalent to
the two equations (V1)
\(\partial_v\mathcal P^{\rm V}_{\rm LG}=(\mathcal M_{\rm LG}/r-\tfrac12)\,\partial_r\mathcal P^{\rm V}_{\rm LG}-\mathcal M_{\rm LG}/(2r^2)\,\mathcal P^{\rm V}_{\rm LG}-\big(i\lambda_{\rm LG}/(2r)+m_{\rm LL}/2\big)\mathcal Q^{\rm V}_{\rm LG}\)
and (V2)
\(\partial_r\mathcal Q^{\rm V}_{\rm LG}=(i\lambda_{\rm LG}/r-m_{\rm LL})\,\mathcal P^{\rm V}_{\rm LG}\),
for every function \(\mathcal M_{\rm LG}(v)\), on both sides of the
apparent horizon. This was derived by computer algebra and is checked
three times: numerically at 80 random points, half of them inside the
apparent horizon, with the spin connection of the library of Part III
(residual \(1.1\times10^{-14}\), check LV-1), and exactly, with the spin
connection computed from the coframe for a symbolic
\(\mathcal M_{\rm LG}(v)\), in Python (YV-1) and in the Wolfram Language
(WV-1); in both, the same spinor with a wrong coefficient in (V1) leaves
a residual, so the checks can fail. The rate \(d\mathcal M_{\rm LG}/dv\)
does not appear in (V1), (V2). Written as
\(\partial_v\mathcal P^{\rm V}_{\rm LG}+\tfrac12 f_{\rm LG}\,\partial_r\mathcal P^{\rm V}_{\rm LG}=\dots\),
(V1) carries \(\mathcal P^{\rm V}_{\rm LG}\) along the outgoing radial
light rays, \(dr/dv=f_{\rm LG}/2\), outward outside the apparent horizon
and inward inside it (WV-5). (V2) contains no derivative in \(v\): each
slice \(v=\text{const}\) is made of ingoing light rays (it is a
\emph{null surface}: its normal vector is light-like), and on it
\(\mathcal Q^{\rm V}_{\rm LG}\) follows from
\(\mathcal P^{\rm V}_{\rm LG}\) by an integral in \(r\). \emph{The
charge.} For every solution,
\(\partial_v\big(2\lvert\mathcal P^{\rm V}_{\rm LG}\rvert^2\big)+\partial_r\big(f_{\rm LG}\lvert\mathcal P^{\rm V}_{\rm LG}\rvert^2-\lvert\mathcal Q^{\rm V}_{\rm LG}\rvert^2\big)=0\)
exactly (YV-2, WV-2). The charge of the brochon is the number of
particles minus the number of antiparticles, the integral of the time
component of its conserved current; on the slice its density, integrated
over the angles, is \(2\lvert\mathcal P^{\rm V}_{\rm LG}\rvert^2\), and
\(f_{\rm LG}\lvert\mathcal P^{\rm V}_{\rm LG}\rvert^2-\lvert\mathcal Q^{\rm V}_{\rm LG}\rvert^2\)
is its flux through the sphere of radius \(r\). The charge
\(N^{\rm V}_{\rm LG}(v)\) on a slice therefore changes only by the
fluxes through its inner and outer ends. \emph{The static limit.} For a
constant mass (\(\mathcal M_{\rm LG}=1\)),
\(\mathcal P^{\rm V}_{\rm LG}=e^{-i\omega^{\rm S}_{\rm LG}v}p^{\rm V}_{\rm LG}(r)\)
and
\(\mathcal Q^{\rm V}_{\rm LG}=e^{-i\omega^{\rm S}_{\rm LG}v}q^{\rm V}_{\rm LG}(r)\)
turn (V1), (V2) into
\(r(r-2)\,p^{\rm V\prime}_{\rm LG}=2i\omega^{\rm S}_{\rm LG}r^2p^{\rm V}_{\rm LG}-p^{\rm V}_{\rm LG}-(i\lambda_{\rm LG}r+m_{\rm LL}r^2)\,q^{\rm V}_{\rm LG}\)
and
\(r\,q^{\rm V\prime}_{\rm LG}=(i\lambda_{\rm LG}-m_{\rm LL}r)\,p^{\rm V}_{\rm LG}\)
(WV-3), with the exponents \(0\) and
\(4i\omega^{\rm S}_{\rm LG}-\tfrac12\) at the horizon, as in stage 5
(YV-3, WV-3); the solution with the exponent 0, analytic at the horizon,
is the one that falls in. These solutions are those of stage 5 in the
freely falling tetrad:
\(p^{\rm V}_{\rm LG}=(1+w_{\rm LG})^{-1/2}e^{i\omega^{\rm S}_{\rm LG}\Delta_{\rm LG}}\mathcal P_{\rm LG}\)
and
\(q^{\rm V}_{\rm LG}=(1+w_{\rm LG})^{1/2}e^{i\omega^{\rm S}_{\rm LG}\Delta_{\rm LG}}\mathcal Q_{\rm LG}\),
where \(\Delta_{\rm LG}(r)=v-\tau_{\rm LG}\) has the derivative
\(1/(1+w_{\rm LG})\) (YV-4, WV-4; the opposite powers of
\(1+w_{\rm LG}\) fail). The reason: the coframe of D10 with a constant
mass is the freely falling coframe of stage 5 boosted in the radial
direction (a Lorentz transformation of the frame, with the velocity
\(\tanh\chi_{\rm LG}\)) with the rapidity
\(\chi_{\rm LG}=-\ln(1+w_{\rm LG})\) (WV-4b). The quasi-bound states of
stage 5 are therefore found again: \(1S_{1/2}\), \(2P_{1/2}\) and
\(2P_{3/2}\) at \(\mathcal M_{\rm LG}m_{\rm LL}=0.3\) to
\(2.1\times10^{-13}\) of \(m_{\rm LL}\) (LV-2), and \(1S_{1/2}\) in
Python, with its own start and Muller's method of finding roots, to
\(8.7\times10^{-14}\) (YV-5) \\
E16 & 22 Nov 2026, 10:22 to 11:39 & stage 6, the evolution in time: the
method and its tests (E15), computed with checks; no manuscript text
(C3). \emph{The method of lines:} \(\mathcal P^{\rm V}_{\rm LG}\) is
represented by its values on a uniform grid of radii from \(r=1.2\),
inside the apparent horizon at all times, to \(r=900\); the derivative
in \(r\) by finite differences of the sixth order (seven points),
\(\mathcal Q^{\rm V}_{\rm LG}\) by the integral of (V2) from the outer
end, where \(\mathcal Q^{\rm V}_{\rm LG}=0\), and a small term
proportional to the eighth difference of the grid values (the
dissipation of Kreiss and Oliger), which damps the shortest wavelengths
that the grid can carry and changes smooth solutions only at the order
\(\delta r_{\rm LG}^7\) in the spacing \(\delta r_{\rm LG}\); the
resulting system of ordinary differential equations in \(v\) is
integrated by CVODE. At the inner end no condition is needed, because
there the light rays along which (V1) carries the field move inward, out
of the grid (E15): the interior of the black hole is cut out of the grid
(\emph{excision}), and what falls in leaves the grid as it leaves the
outside world. The outgoing rays move with
\(dr/dv=f_{\rm LG}/2<\tfrac12\), so that nothing from inside \(r=225\)
reaches the outer end during the runs. \emph{Tests.} For a constant mass
the state \(1S_{1/2}\) at \(\mathcal M_{\rm LG}m_{\rm LL}=0.3\) evolves
as \(e^{-i\omega^{\rm S}_{\rm LG}v}\) times its initial profile, its
decay by absorption included, to \(4.7\times10^{-9}\) for
\(\delta r_{\rm LG}=0.1\) and \(4.5\times10^{-11}\) for \(0.05\)
(relative to the largest value; LV-3; the ratio of the two, about
\(2^{6.7}\), agrees with the order 6 of the scheme). The change of the
charge equals the integrated flux through the two ends to
\(2.1\times10^{-8}\) of the charge (LV-4), and to \(2.6\times10^{-8}\)
in the four runs below (LV-5). \\
E17 & 22 Nov 2026, 10:22 to 11:39 & stage 6, the runs with a changing
mass (E16), computed with checks; no manuscript text (C3). \emph{The
runs} (mass of the black hole 1 at the start, \(m_{\rm LL}=0.2\),
\(\delta r_{\rm LG}=0.1\)): the brochon starts in the state \(1S_{1/2}\)
(\(y_{\rm LG}=0.2\)). From \(v=50\) the mass changes smoothly (a step
whose derivatives of every order are continuous) to \(1.25\)
(\emph{growth}, \(y_{\rm LG}=0.25\) at the end) or to \(0.8\)
(\emph{shrinking}, D10; \(y_{\rm LG}=0.16\)), within the time 1000
(\emph{slow}) or 20 (\emph{fast}); then it stays constant for the time
300. Two limits serve as references. \emph{Adiabatic}: a change so slow
that the brochon stays in the state \(1S_{1/2}\) of the instantaneous
mass (the adiabatic theorem of quantum mechanics, for slow changes of a
system with discrete levels); its charge would then decay as
\(\exp\int2\,{\rm Im}\,\omega^{\rm S}_{\rm LG}\,dv\) with the
instantaneous frequency of \(1S_{1/2}\). \emph{Sudden}: a change so fast
that the field has no time to change, so that it is the old state
\(1S_{1/2}\) written in terms of the states of the new mass. All charges
below are relative to the initial charge. Growth, slow: at the end of
the change the charge is \(0.0386\), the adiabatic charge \(0.0332\),
and the charge of the state \(1S_{1/2}\) (from the decomposition of E18)
\(0.0332\), which is \(0.999\) of the adiabatic charge: the brochon
follows the growing mass adiabatically; the state \(2S_{1/2}\) holds
\(0.0040\). Growth, fast: the change is sudden; the projection of the
field on the state \(1S_{1/2}\) of the final mass at the end of the
change, \(0.9407\), is close to that of the old state, \(0.9360\).
Shrinking, slow: at the end of the change the charge of \(1S_{1/2}\) is
\(0.921\) of the adiabatic charge, and the total charge \(0.3794\) lies
below the adiabatic \(0.4001\): part of the field left the state
\(1S_{1/2}\) and was absorbed. That the slow shrinking is less adiabatic
than the slow growth agrees with the adiabatic theorem: its time scale
is set by the distance of the levels, which falls with \(y_{\rm LG}^2\),
so that the same duration of the change is less slow at the smaller
mass. Shrinking, fast: sudden; projection \(0.9345\), old state
\(0.9333\). At the end of the runs the growth has made the absorption
faster (the decay rate of \(1S_{1/2}\) rises steeply with
\(y_{\rm LG}\), E13) and the shrinking slower: the final charges are
\(7.388\times10^{-3}\), \(0.2561\), \(0.3505\) and \(0.7994\) (growth
slow and fast, shrinking slow and fast), against the adiabatic
\(6.500\times10^{-3}\), \(0.1700\), \(0.3643\) and \(0.8306\). After a
growth more charge remains than in the adiabatic limit (\(1.137\) and
\(1.507\) times), because the parts of the field outside \(1S_{1/2}\)
are absorbed more slowly than \(1S_{1/2}\) at \(y_{\rm LG}=0.25\); after
a shrinking less (\(0.962\) and \(0.962\) times). For comparison: a
hydrogen atom whose nuclear charge is multiplied suddenly by 1.25 or by
0.8 keeps its ground state with the probability
\([4ZZ'/(Z+Z')^2]^3=0.963\) (non-relativistic, without absorption); the
computed charge of \(1S_{1/2}\) in the old state is \(0.9535\) for the
growth and \(0.9329\) for the shrinking. \\
E18 & 22 Nov 2026, 10:22 to 11:39 & stage 6, the meaning of the
projections and the decomposition of the field into states (E17). The
projection, or \emph{occupation}, of the field
\(\mathcal P^{\rm V}_{\rm LG}\) on a state \(p_n\) is
\(O_{n,{\rm LG}}=\lvert\langle p_n,\mathcal P^{\rm V}_{\rm LG}\rangle_{\rm LG}\rvert^2/(\langle p_n,p_n\rangle_{\rm LG}\langle\mathcal P^{\rm V}_{\rm LG},\mathcal P^{\rm V}_{\rm LG}\rangle_{\rm LG})\),
with the charge product
\(\langle a,b\rangle_{\rm LG}=\int2\bar a\,b\,dr\) on the slice. The
quasi-bound states lose charge into the black hole, the evolution does
not conserve the charge product, and the states are not orthogonal: the
normalized overlap
\(\lvert\langle p_1,p_2\rangle_{\rm LG}\rvert/(\langle p_1,p_1\rangle_{\rm LG}\langle p_2,p_2\rangle_{\rm LG})^{1/2}\)
of \(1S_{1/2}\) and \(2S_{1/2}\) is \(0.135\) at \(y_{\rm LG}=0.2\),
\(0.292\) at \(0.25\) and \(0.048\) at \(0.16\) (zero for orthogonal
states). The occupations are therefore not probabilities: at the end of
the fast growth those of \(1S_{1/2}\) and \(2S_{1/2}\) add up to
\(1.1609\). \emph{The decomposition} writes the field as
\(\sum_nc_{n,{\rm LG}}p_n\) over the states \(1S_{1/2}\) to \(4S_{1/2}\)
of the final mass (frequencies at \(y_{\rm LG}=0.25\): \(1S_{1/2}\):
\(0.958838-0.0136\,i\); \(2S_{1/2}\): \(0.990251-1.92\times10^{-3}\,i\);
\(3S_{1/2}\): \(0.995893-5.40\times10^{-4}\,i\); \(4S_{1/2}\):
\(0.997769-2.18\times10^{-4}\,i\); at \(0.16\): \(1S_{1/2}\):
\(0.984759-7.80\times10^{-4}\,i\); \(2S_{1/2}\):
\(0.996401-1.00\times10^{-4}\,i\); \(3S_{1/2}\):
\(0.998453-2.87\times10^{-5}\,i\); \(4S_{1/2}\):
\(0.999146-1.19\times10^{-5}\,i\), all as
\(\omega^{\rm S}_{\rm LG}/m_{\rm LL}\)), with the coefficients of the
best approximation in the charge (least squares: the linear equations
whose matrix, the \emph{Gram matrix}, holds the charge products of the
states), plus a remainder;
\(W_{n,{\rm LG}}=\lvert c_{n,{\rm LG}}\rvert^2\langle p_n,p_n\rangle_{\rm LG}\)
is the charge that the state \(n\) would carry alone. Because of the
overlaps, the \(W_{n,{\rm LG}}\) do not add up to the total charge
either, and the result depends on how the coefficients are fitted. Two
fits were made: at the single time of the end of the change, and over
the whole time after the change, where every coefficient must evolve as
\(e^{-i\omega^{\rm S}_{\rm LG}v}\) with the frequency of its state. For
the slow growth and both shrinking runs they agree within 2 \% (the
charge of \(1S_{1/2}\) at the end of the change: \(0.0332\) and
\(0.0332\), \(0.3686\) and \(0.3709\), \(0.8462\) and \(0.8322\)); for
the fast growth they give \(0.8476\) and \(0.9494\). The coefficient of
\(1S_{1/2}\) found at a single time at the end of the change and at the
end of the run follows \(e^{-i\omega^{\rm S}_{\rm LG}v}\) to
\(6.5\times10^{-3}\) and \(0.02\) after the shrinking, but only to
\(0.053\) and \(0.13\) after the growth. The reason: at
\(y_{\rm LG}=0.25\) the decay rate
\(\lvert{\rm Im}\,\omega^{\rm S}_{\rm LG}\rvert\) of \(1S_{1/2}\) is
\(0.433\) of the distance
\({\rm Re}\,\omega^{\rm S}_{\rm LG}(2S_{1/2})-{\rm Re}\,\omega^{\rm S}_{\rm LG}(1S_{1/2})\)
of the two levels (at \(0.16\): \(0.067\)): the levels are nearly as
broad as they are apart, and the charge cannot be divided among the
states better than about 10 \%. The remainder outside the four states
holds 1 \% to 4 \% of the charge in every run (\(0.0195\) at the end of
the fast growth, \(0.0436\) at the end of the fast shrinking). The
decline of the occupation of \(1S_{1/2}\) during the slow growth
(\(0.9056\) at the end of the change, \(0.6069\) at the end) is thus not
a loss from \(1S_{1/2}\): that state keeps its adiabatic charge, but it
is absorbed faster than the admixture of \(2S_{1/2}\) that the change
produced, whose share of the remaining field grows \\
E19 & 22 Nov 2026, 10:22 to 11:39 & the programs of E15 to E18 and the
corrections on the way. Rust: the crate \texttt{rust/lgvaidya}, program
\texttt{lg\_vaidya} (ALL PASS (8 of 8); about 5 minutes), with the
separation (LV-1), the static limit (LV-2), the evolution (LV-3, LV-4),
the four runs (LV-5), the grid (LV-6: the fast runs repeated with
\(\delta r_{\rm LG}=0.2\) agree with \(0.1\) to \(5.4\times10^{-8}\) in
the final occupation, the final charge and the charges of \(1S_{1/2}\)
and \(2S_{1/2}\)), the decomposition (LV-7: the initial state decomposed
in the states of the initial mass gives \((1,0,0,0)\) to
\(1.7\times10^{-18}\)) and the identification of the states (LV-8:
\(1S_{1/2}\) and \(2S_{1/2}\) of the final mass of the growth, followed
from \(y_{\rm LG}=0.2\) in steps of 1 \% of the mass, equal those of
stage 5 at \(y_{\rm LG}=0.25\) to \(6.3\times10^{-14}\); every step of
every state is guarded: the root must lie nearer to the small-mass level
\(1-y_{\rm LG}^2/(2n^2)+F^{\rm H}_{\rm LG}y_{\rm LG}^4\) of its \(n\)
than half the distance to the neighbouring levels, or the program
stops). Data: \texttt{data/lg\_vaidya\_values.csv} (the numbers quoted
here), \texttt{data/lg\_vaidya\_runs.csv} (the four runs at every output
time). Python: \texttt{python/lg\_vaidya.py} (ALL PASS (9 of 9); 30 s):
YV-1 to YV-4 by sympy (E15); YV-5, the state \(1S_{1/2}\) at
\(y_{\rm LG}=0.3\); YV-6 and YV-7, the fast growth recomputed to
\(v=120\) with its own operators (weights of the finite differences and
of the integrals from the equations of their moments, solved by numpy)
and CVODE: the charges agree with the Rust program to
\(6.1\times10^{-14}\), the balance holds to \(6.0\times10^{-9}\); YV-8,
at \(v=120\) the state \(1S_{1/2}\) of the mass 1.25 equals that of
stage 5 and the occupations agree with the Rust program
(\(3.0\times10^{-13}\)); YV-9, the decomposition at the end of the
change agrees with it (\(1.5\times10^{-12}\)). Wolfram:
\texttt{wolfram/source/lg\_vaidya.wls} (ALL PASS (6 of 6), exact;
request 342, 1.5 s): WV-1 to WV-5 (E15). Fragments:
\texttt{code/scripts/write\_lg\_vaidya\_fragments.py}. Corrections
before any result was used: (1) the first version of the evolution
reached the right-hand side through a raw pointer; it now uses an owned
copy of its data. (2) The first runs stopped (CVODE reached its maximum
number of steps near \(r=10^7\)) when the search for the state
\(1S_{1/2}\) of a changed mass, started from the small-mass formula,
converged to a spurious root and the profile was integrated far outside;
the states are now followed in the mass in steps of 1 \%, with
assertions. (3) The start from the small-mass formula failed at
\(y_{\rm LG}=0.25\); the state is now followed from \(y_{\rm LG}=0.2\).
(4) LV-5 failed on the first production grid, \(\delta r_{\rm LG}=0.2\)
(balance \(1.25\times10^{-6}\), tolerance \(10^{-6}\)); the production
grid became \(\delta r_{\rm LG}=0.1\), and LV-6 compares the two. (5) A
first version of LV-8 identified the states by the number of sign
changes of the real part of their profiles and failed: on a slice
\(v=\text{const}\) a profile winds as
\(e^{i\omega^{\rm S}_{\rm LG}r_*}\) with the tortoise coordinate
\(r_*=r+2\mathcal M_{\rm LG}\ln(r/2\mathcal M_{\rm LG}-1)\), which gave
18 to 59 sign changes instead of 0 to 3; with the winding removed the
counts are \(n-1\) at \(y_{\rm LG}=0.16\), but not for \(1S_{1/2}\) and
\(2S_{1/2}\) at \(0.25\), whose profiles are then far from real. The
comparison with stage 5 replaced it. (6) An edit of the program deleted
seven quantities per run from \texttt{data/lg\_vaidya\_values.csv}; the
fragment writer stopped on the first missing one, the lines were
restored and the program rerun, and every other number came out the same
(the file of the runs byte for byte). (7) A trial computation in Python,
not kept, followed the states by Muller's method from the small-mass
formula without the guard and converged twice to other states; its
decompositions were meaningless and were not used \\
E20 & 22 Nov 2026, 11:39 to 11:54 & stage 7a, the pair of D11, computed
with checks; no manuscript text (C3). \emph{The levels.} The pair is
hydrogen-like: with the reduced mass
\(\mu^{\rm pair}_{\rm LG}=m^{\rm pair}_{\rm LG}/2\) of two particles of
mass \(m^{\rm pair}_{\rm LG}\), its levels are
\(2m^{\rm pair}_{\rm LG}-m^{\rm pair}_{\rm LG}(\alpha^{\rm Z}_{\rm LG})^2/(4n^2)\)
to the leading order, its size is the Bohr radius
\(2/(m^{\rm pair}_{\rm LG}\alpha^{\rm Z}_{\rm LG})\), and
\(\alpha^{\rm Z}_{\rm LG}=\) \(0.00975\). The next order,
\(m^{\rm pair}_{\rm LG}(\alpha^{\rm Z}_{\rm LG})^4\), follows from the
\emph{Breit-Fermi Hamiltonian}: the Hamiltonian of two Dirac particles
that exchange one massless vector boson, expanded to the second order in
their velocities, which for the electron and the positron gives the
levels of positronium. Its terms are the relativistic correction of the
kinetic energy, the Darwin term (a contact term of the zitterbewegung),
the orbit-orbit term (the magnetic interaction of the two currents and
the retardation of the exchange), the spin-orbit term, and the
interaction of the two magnetic moments (a contact term and a tensor
term). They split the levels of the same \(n\) into the \emph{fine
structure} (by the orbital angular momentum \(\ell\) and the total
angular momentum \(j\)) and the \emph{hyperfine structure} (by the total
spin \(s\): \(s=0\), the singlet, \emph{para}, and \(s=1\), the triplet,
\emph{ortho}); the states are written \(n\,{}^{2s+1}L_j\) with
\(L=S,P,D\) for \(\ell=0,1,2\). In first-order perturbation theory in
the states \(|n\,\ell\,s\,j\rangle\) the coefficients of
\(m^{\rm pair}_{\rm LG}(\alpha^{\rm Z}_{\rm LG})^4\) are exact rational
numbers; without the annihilation term, for \(n\le2\): \(1{}^{1}S_{0}\):
\(-\tfrac{21}{64}\); \(1{}^{3}S_{1}\): \(\tfrac{1}{192}\);
\(2{}^{1}S_{0}\): \(-\tfrac{53}{1024}\); \(2{}^{3}S_{1}\):
\(-\tfrac{31}{3072}\); \(2{}^{1}P_{1}\): \(-\tfrac{31}{3072}\);
\(2{}^{3}P_{2}\): \(-\tfrac{43}{15360}\); \(2{}^{3}P_{1}\):
\(-\tfrac{47}{3072}\); \(2{}^{3}P_{0}\): \(-\tfrac{95}{3072}\), all with
the leading terms \(-\tfrac14\) and \(-\tfrac1{16}\) of
\(m^{\rm pair}_{\rm LG}(\alpha^{\rm Z}_{\rm LG})^2\) removed (the table
for \(n\le4\) is \texttt{data/lg\_pair\_levels.csv}). With the term of
virtual annihilation into one zeugon, which acts only on the triplet S
states, \(1\,{}^3S_1\) becomes \(\tfrac{49}{192}\). The same operators
give exactly Dirac's fine structure for one particle in a fixed Coulomb
field (YP-1) and Ferrell's formula for positronium (YP-2; S. R. Ferrell,
1951, as quoted in arXiv:1712.02022, equations (9) to (11)). \emph{The
numbers} (units of the cutoff \(\Lambda_{\rm LL}\); quasi-particles of
mass \(E^{\rm CJT,gap}_{\rm SF}=\) \(0.3664\), with the mass
\(m_{\rm LL}=0.2\) in brackets): Bohr radius \(560\) (\(1026\)) in units
of \(1/\Lambda_{\rm LL}\); binding energy of the ground state
\(8.70\times10^{-6}\) (\(4.75\times10^{-6}\)), \(2.38\times10^{-5}\) of
the mass of one quasi-particle; hyperfine splitting of the ground state
\(1.10\times10^{-9}\) without and \(1.93\times10^{-9}\) with
annihilation; spread of the levels \(n=2\) \(1.62\times10^{-10}\). If
the pair can annihilate as positronium does, the formulas of positronium
give the decay rates \(1.61\times10^{-11}\) of \(1\,{}^1S_0\) into two
zeugons and \(1.93\times10^{-14}\) of \(1\,{}^3S_1\) into three (D11
(i)). \emph{Tidal dissociation.} In the frame of an observer who falls
freely at the radius \(r\) of a Schwarzschild black hole of
gravitational radius \(\mathcal M_{\rm LG}=G_{\rm N}M_{\rm LG}/c^2\),
the black hole stretches the pair along the radius and squeezes it
across: the relative coordinate feels the tidal potential
\(-k^{\rm tid}_{\rm LG}z^2+\tfrac12k^{\rm tid}_{\rm LG}(x^2+y^2)\) with
\(k^{\rm tid}_{\rm LG}=\mu^{\rm pair}_{\rm LG}\mathcal M_{\rm LG}/r^3\)
(from the deviation of neighbouring geodesics). On the radial axis the
barrier of \(-\alpha^{\rm Z}_{\rm LG}/z-k^{\rm tid}_{\rm LG}z^2\) has
the top
\(-\tfrac32(\alpha^{\rm Z}_{\rm LG})^{2/3}(2k^{\rm tid}_{\rm LG})^{1/3}\);
when it falls below the level \(n\), the pair is freed classically, at
\((r^{\rm ion}_{\rm LG})^3=54n^6\mathcal M_{\rm LG}/((\mu^{\rm pair}_{\rm LG})^2(\alpha^{\rm Z}_{\rm LG})^4)\).
This lies outside the horizon if
\(\mathcal M_{\rm LG}<\tfrac{3\sqrt3}2n^3/(\mu^{\rm pair}_{\rm LG}(\alpha^{\rm Z}_{\rm LG})^2)\):
for the ground state of the pair of quasi-particles,
\(\mathcal M_{\rm LG}<\) \(1.492\times10^{5}\)\(/\Lambda_{\rm LL}\),
i.e.~black holes lighter than \(1.492\times10^{5}\)
\(M_{\rm Pl}^2/\Lambda_{\rm LL}\) (\(M^2_{\rm Pl}=\hbar c/G_{\rm N}\))
tear the pair apart before it reaches the horizon, heavier ones swallow
it whole. The estimate assumes a static tidal field and the unperturbed
level (no tunnelling, no shift of the level by the tidal field); it
gives the scale. Since \(\Lambda_{\rm LL}\) is not known in physical
units (D11 (iii)), only illustrative values can be given: the critical
mass is \(3.97\times10^{25}\) kg (\(1.99\times10^{-5}\) solar masses)
for \(\Lambda_{\rm LL}=1\) eV, \(3.97\times10^{16}\) kg for 1 GeV and
\(3.97\times10^{13}\) kg for 1 TeV; for any \(\Lambda_{\rm LL}\) above
about 1 eV, black holes of stellar mass swallow the pairs whole \\
E21 & 22 Nov 2026, 11:39 to 11:54 & the programs of E20. Python:
\texttt{python/lg\_pair.py} (ALL PASS (3 of 3); exact, sympy; 16 s):
YP-1 Dirac's fine structure, YP-2 Ferrell's formula, YP-3 the radial
integrals, each for \(n\le5\); it writes
\texttt{data/lg\_pair\_levels.csv}. Rust: the crate
\texttt{rust/lgpair}, program \texttt{lg\_pair} (ALL PASS (2 of 2);
under 1 s): LP-1 the coefficients by Gauss-Laguerre quadrature of the
radial integrals (48 points), equal to the exact ones to
\(2.3\times10^{-14}\) for all states with \(n\le4\); LP-2 the tidal
radius found numerically (golden-section search for the top of the
barrier, bisection for the radius) equal to the closed form; it writes
\texttt{data/lg\_pair\_values.csv} (the numbers of E20). Wolfram:
\texttt{wolfram/source/lg\_pair.wls} (ALL PASS (2 of 2); exact): WP-1
the coefficients for \(n\le3\) with the radial functions of the Wolfram
Language, WP-2 the barrier. Fragments:
\texttt{code/scripts/write\_lg\_pair\_fragments.py}. Correction: in the
first run of the Wolfram script (request 343) WP-2 failed: the position
of the top and the radius were found with \texttt{Solve}, whose results
are \texttt{Root} objects that \texttt{Simplify} did not reduce to the
closed forms (the formulas themselves were right: the printed
\texttt{Root} objects are the cube roots of the closed forms); the
script now checks the closed forms directly (request 344) \\
C4 & 22 Nov 2026, 13:46 & the author's answer to my question (D11), sent
after the commit \texttt{e38f21e} (11:54 PST): the predictions that
``can be detected gravitationally'' are to be made with black holes. He
asked which other questions I had \\
C5 & 22 Nov 2026, 13:51 & the author's answer to my further questions
(the value of \(\Lambda_{\rm LL}\) in physical units, which effects of
black holes, rotating black holes, the particles of the pair, the
annihilation of the pair): ``I would rather not decide; take the best
route'' (the route: decision D12); and his instruction to proceed with
all of my manuscripts, less tersely, using the remarks that he provided
to lighten up and to be more user-friendly. The writing of this folder
resumes (Section 4.1, the plan) \\
E22 & 22 Nov 2026, 14:57 to 15:39 & the rewriting of Parts I to III of
\texttt{lord\_gravity.md} after C5 (Section 4.1), in the friendlier
style that the author asked for. \emph{Added:} the paragraph \emph{How
to read this manuscript}; Section 0.1, a word list of about fifty words
in four groups (space-time and gravitation; Lord's construction; the
fields and their charges; the superfluid), each explained in a few
sentences (\texttt{docs\_src/parts/glossary\_lg.part.md}); Section 0.2,
the results in brief; a paragraph \emph{In words} at the head of
Sections 2, 3, 5, 6 and 13 to 16 and of Parts II and III; Appendix A,
the algebra of the Lorentz super-loop with an example of every object
(\texttt{algebra\_lg.part.md}, made from the appendix of the documents
of the folder of the revision for the Higgs field, with the notation of
this folder); Appendix B, the brochon model
(\texttt{brochon\_lg.part.md}, see E23). \emph{Changed:} the 30 lines of
the source that took objects from the earlier hypotheses or from the
earlier work, and one more in which the words were split over two lines,
now name the brochon model and its appendix; the numbers of the ground
state that Section 15.4, Section 17 and a remark typed by hand are now
fragments. The assembler \texttt{code/scripts/assemble\_docs.py} has a
new placeholder, PART, which includes a file of
\texttt{docs\_src/parts/} in full. HCK-H1 and HCK-H2 pass \\
E23 & 22 Nov 2026, 14:57 to 15:39 & Appendix B of E22 and the worked
examples. \emph{Appendix B} states the field and its charges, GEE and
the covariant derivative with the proof of covariance, the Lagrangian
and the field equations, the mean field and the methods HS, KS and CJT
with the theorem KS = CJT, the effective couplings and the stability of
the normal vacuum, the recursion of the interaction (\(V^{(0)}\),
\(V^{(1)}\), \(V^{(2)}\)), the polar ground state with its
quasi-particles and excitations, the breaking to
\(SL(2,\mathbb C)\times O(2)\), the masses of the gauge bosons and the
zeugon; its numbers are copied by
\texttt{code/scripts/write\_lg\_examples\_fragments.py} (now 44
fragments of the programs of Part II of the manuscript of the superfluid
in the folder of the revision for the Higgs field, with the SHA-256
checksums of the twelve files that it reads). \emph{The worked examples}
of Sections 2.4, 11.1, 16.1 and Appendix B (Lord's change of
representation by the boost of Appendix A.3 and its frame, a set of four
Hermitian matrices that is not a representation, a plane wave of the
brochon with its current and energy-momentum tensor, the quasi-particles
and the gap equation of the polar state) are computed by the new Rust
program \texttt{lg\_worked} (crate \texttt{rust/lgworked}; VERDICT: ALL
PASS (30 of 30)) and checked by \texttt{python/lg\_worked.py} (VERDICT:
ALL PASS (28 of 28); exact) and \texttt{wolfram/source/lg\_worked.wls}
(VERDICT: ALL PASS (5 of 5); requests 345 to 347); fragments by
\texttt{code/scripts/write\_lg\_worked\_fragments.py} \\
E24 & 22 Nov 2026, 14:57 to 15:39 & corrections made while E22 and E23
were written, before the commit. (1) The Wolfram script
\texttt{lg\_worked.wls}: in request 345 the check WW-5 passed with a
warning of the precision of \texttt{N}, which had evaluated an
expression that is exactly zero; in request 346 it compared the result
of \texttt{Integrate}, written with \texttt{ArcCoth}, with the logarithm
of (A.19) by \texttt{Simplify}, which failed (WW-5 FAIL); request 347
proves the identity through the hyperbolic cotangent of the logarithm
and compares the numbers side by side: 5 of 5, no warning. (2) A draft
of Appendix B compared the strength of the zeugon potential with the
fine-structure constant as ``about three quarters'' of it; the ratio is
about 1.34; the comparison was removed. (3) A fragment that contains
words was printed inside a formula (``1.113to1.184''); it is now printed
as text. (4) The list of the files read by
\texttt{write\_lg\_examples\_fragments.py} with their checksums
overflowed the page; it is now a block of code. (5) A replacement string
of an editing script was written without the raw prefix of Python, so
that the characters of \texttt{\textbackslash{}nu} became a line break;
the script stopped at its own check before it wrote anything \\
E25 & 22 Nov 2026, 15:35 to 15:39 & the times of the log checked against
the times of the commits: several entries ended after the commits that
contain them, and E20 and E21 lay entirely after theirs. Recorded as a
false statement in Section 4.3 and corrected there and in the entries
E1, E3 to E6, E11 to E21 and C4, and in the record of Section 4.2 \\
E26 & 22 Nov 2026, 15:47 to 16:17 & Part IV of \texttt{lord\_gravity.md}
written from stage 5 (E7 to E14), in the style of C5: Sections 18 to 22
(the particle, the tetrads and the separation; the horizon and the
inside; the quasi-bound states; the greybody factors; what Part IV
establishes), each with a paragraph \emph{In words}, every term defined
at its first use (separation of variables, partial wave, regular
singular point and exponents, analytic solution, quasi-bound state,
Wronskian, secant method, fine structure, greybody factor), worked
examples (the angular functions of \(j=\tfrac12\) checked by hand in the
text; the state \(1S_{1/2}\) at \(\mathcal M_{\rm LG}m_{\rm LL}=0.1\);
the small-mass formulas evaluated for \(1S_{1/2}\)), a table of the six
states at two masses (Table IV.1) and two figures (Figure IV.1, the
solution on both sides of the horizon; Figure IV.2, the decay rates
against the small-mass formulas), drawn by the new program
\texttt{python/lg\_schw\_figures.py} from the data of \texttt{lg\_schw}
(the palette of the other figures of the project, checked for
colour-blind readers). The word list has a new group, \emph{Waves around
a black hole}. The provenance and the symbols of the main text moved
from Sections 18 and 19 to Appendices C and D, so that later parts do
not renumber them; every reference was changed.
\texttt{code/scripts/write\_lg\_schw\_fragments.py} writes 17 more
fragments (the table; the deviations of LS-2, LS-4, LS-5, LS-7, LS-9a,
LS-9b, LS-14a and LS-14c and the smallest residual of LS-14b; the weight
of the second part of the quasi-particle of Theorem L3; the frequency of
\(1S_{1/2}\) at \(0.1\) and its distance below the rest energy; the
decay times in periods; the gravitational radius of the Sun and the
particle mass at which a black hole of one solar mass has
\(\mathcal M_{\rm LG}m_{\rm LL}=0.1\), from the nominal solar mass
parameter of the IAU and \(\hbar c\) of CODATA 2018). Corrections made
while writing, before the commit: the draft said that \(3D_{3/2}\) at
\(\mathcal M_{\rm LG}m_{\rm LL}=0.01\) lives ``\(10^{27}\) periods''
(the computed number is \(5\times10^{25}\)) and that \(1S_{1/2}\) at
\(0.1\) lives ``\(1.8\times10^4\) periods'' (\(5.7\times10^{3}\)); and
that every particle heavier than \(1.3\times10^{-11}\) eV ``is absorbed
within a few periods'' around a black hole of a solar mass, which the
computations do not show (at large \(\mathcal M_{\rm LG}m_{\rm LL}\) a
particle can orbit like a classical body); the text now states only what
was computed. A first version of an added Rust string lost its line
continuations to the Python editing script (non-raw string, E24 (5));
the report line was rewrapped before the commit \\
E27 & 22 Nov 2026, 15:47 to 16:17 & Theorem L3 (Section 18.1 of the
manuscript), found while E26 was written: if \(\chi\) solves the Dirac
equation with the mass \(E^{\rm CJT,gap}_{\rm SF}\), then
\(\chi-D^{\rm CJT}_{\rm SF}/(E^{\rm CJT,gap}_{\rm SF}+M^{\rm CJT,D}_{\rm SF})\,\mathcal C_{\rm LG}\chi\)
solves the equation (15.11) of the quasi-particle, and \(\chi\) and
\(\mathcal C_{\rm LG}\chi\) have the same currents and scalar densities.
D9 and Section 15.4 had rested on Theorem L2 alone, which gives the
other direction (every solution of (15.11) solves the equation of the
second order of the Dirac particle); the use of the Dirac solutions for
the quasi-particle needs this direction, which was not shown until now.
The statement of D9 stands; its reason is complete, and Section 15.4 now
names Theorem L3. Checks: Rust LS-14a (the separated spinor with the
mass of the gap and complex frequencies at 80 random points on both
sides of the horizon, residual \(7.0\times10^{-15}\)), LS-14b (with the
wrong sign of the coefficient the residual is at least \(0.66\): the
check can fail), LS-14c (equal currents and scalar densities,
\(1.2\times10^{-15}\)); Python YS-13 (the same, exactly, by sympy, with
the matrix identities). The new checks use random numbers of their own;
the data of the checks LS-1 to LS-13 are unchanged (compared byte by
byte). Also added: YS-14, which proves that the radial current
\(J_{\rm LG}\) times the angular integral is the flux of the conserved
current into the sphere of radius \(r\), and that the density is
\((\lvert\mathcal P_{\rm LG}\rvert^2+\lvert\mathcal Q_{\rm LG}\rvert^2)\)
times it, as the earlier entries and Part IV state. No Wolfram check was
added: Theorem L3 follows by algebra from the properties of
\(\mathcal C_{\rm LG}\) that the Wolfram script of Part III checks
exactly (WG-6), and the two checks above test it at the solutions.
Verdicts now: Rust ALL PASS (21 of 21), Python ALL PASS (14 of 14),
Wolfram ALL PASS (9 of 9) \\
\end{longtable}

\hypertarget{the-symbols-of-this-document}{%
\section{8. The symbols of this
document}\label{the-symbols-of-this-document}}

\begin{longtable}[]{@{}
  >{\raggedright\arraybackslash}p{(\columnwidth - 2\tabcolsep) * \real{0.5000}}
  >{\raggedright\arraybackslash}p{(\columnwidth - 2\tabcolsep) * \real{0.5000}}@{}}
\toprule\noalign{}
\begin{minipage}[b]{\linewidth}\raggedright
symbol
\end{minipage} & \begin{minipage}[b]{\linewidth}\raggedright
meaning
\end{minipage} \\
\midrule\noalign{}
\endhead
\bottomrule\noalign{}
\endlastfoot
M31 & the manifold of the Hermitian bi-quaternions: Minkowski
space-time \\
\(\eta^{\rm LL}\) & the Minkowski metric \(\mathrm{diag}(1,-1,-1,-1)\)
of M31 \\
\(SL(2,\mathbb C)\) & the group of the complex \(2\times2\) matrices of
determinant 1: the unit bi-quaternions, a subgroup of the Lorentz
super-loop \\
\(\Psi_{\rm LL}\) & the brochon: the isodoublet of Dirac spinor fields
of the earlier hypotheses \\
\(m_{\rm LL}\) & the mass of the brochon \\
GEE & the gauge group
\(SL(2,\mathbb C)\times SU(2)_{\rm v}\times U(1)^{\rm F}_{\rm LL}\) of
the earlier hypotheses \\
\(SU(2)_{\rm v}\) & the isospin group of the earlier hypotheses \\
\(U(1)^{\rm F}_{\rm LL}\) & the group of the phases of the fermion
number of the brochon \\
\(g^{\rm v}_{\rm LL}\), \(g^{\rm F}_{\rm LL}\) & the couplings of
\(SU(2)_{\rm v}\) and of \(U(1)^{\rm F}_{\rm LL}\) \\
\(T^3_{\rm LL}\), \(Q_{\rm LL}\) & the third isospin generator, and the
isospin charge \(Q_{\rm LL}=T^3_{\rm LL}\) (\(\pm\tfrac12\) on the
brochon) \\
\(O(2)\) & the group of the rotations about the direction of the
condensate and the reflection that reverses it \\
LG & the suffix of the symbols that this folder introduces \\
\(\gamma^a_{\rm LL}\), \(\gamma^{ab}_{\rm LL}\) & the Dirac matrices of
the chiral basis of the project, with a frame index, and their
antisymmetrized products \\
\(G_{\rm N}\) & Newton's constant \\
\(G^{\rm LG}_{\mu\nu}\), \(T^{\rm LG}_{\mu\nu}\) & the Einstein tensor
and the energy-momentum tensor of the sources \\
\(\omega^{\rm LG}_{\mu\,ab}\) & the spin connection of a tetrad, defined
by the tetrad postulate \\
\(\Gamma^{\rm LL}_\mu\) & the spin connection of GEE, the gauge field of
its \(SL(2,\mathbb C)\) factor \\
\(R^{\rm LG}_{ab\mu\nu}\), \(R^{\rm LG}\) & the Riemann tensor (two
indices in the frame) and the scalar curvature \\
\(h_{\rm LG}(r)\), \(r\), \(\theta\) & an arbitrary function in a
static, spherically symmetric coframe; the radius and the polar angle \\
\(E^{\rm CJT,gap}_{\rm SF}\) & the gap of the polar superfluid of the
earlier hypotheses (method CJT): the mass of its quasi-particle \\
\(M_{\rm LG}\), \(\mathcal M_{\rm LG}\) & the mass of the black hole,
and \(G_{\rm N}M_{\rm LG}\), the same mass in units of length (half the
radius of the horizon) \\
\(f_{\rm LG}\) & \(1-2\mathcal M_{\rm LG}/r\) \\
\(t\), \(\varphi\) & the time coordinate of Schwarzschild and the
azimuth \\
\(w_{\rm LG}\), \(\tau_{\rm LG}\) & the speed
\((2\mathcal M_{\rm LG}/r)^{1/2}\) of an observer who falls freely from
rest at infinity, and his proper time (the time coordinate of Painlevé
and Gullstrand) \\
\(\omega^{\rm S}_{\rm LG}\) & the frequency of a mode of the brochon in
the Schwarzschild field (complex for a quasi-bound state) \\
\(j_{\rm LG}\), \(m_{j,{\rm LG}}\) & the total angular momentum of a
mode and its component along the axis \\
\(\lambda_{\rm LG}\) & the separation constant of the angular equations,
\(\pm(j_{\rm LG}+\tfrac12)\) \\
\(\kappa^{\rm H}_{\rm LG}\) & the number of Dirac's theory of the
hydrogen atom that tells apart the two states of the same
\(j_{\rm LG}\); \(\kappa^{\rm H}_{\rm LG}=-\lambda_{\rm LG}\) \\
\(n_{\rm LG}\), \(\ell_{\rm LG}\) & the principal quantum number and the
orbital angular momentum of a hydrogen-like level \\
\(\alpha_{\rm LG}(\theta)\), \(\beta_{\rm LG}(\theta)\) & the two
angular functions of a mode \\
\(\mathcal P_{\rm LG}(r)\), \(\mathcal Q_{\rm LG}(r)\) & the two radial
functions of a mode in the freely falling tetrad \\
\(s_{\rm LG}\) & an exponent of the radial equations at a regular
singular point \\
\(J_{\rm LG}\) & the radial current of a mode of real frequency \\
\(\mathcal A_{\rm LG}\) & the greybody factor: the fraction of the flux
of an incoming partial wave that the black hole absorbs \\
\(S_{\rm LG}\) & \(j_{\rm LG}-\ell_{\rm LG}=\pm\tfrac12\), the
projection of the spin of a hydrogen-like level \\
\(N_{\rm LG}\) & the integral of
\(\lvert\mathcal P_{\rm LG}\rvert^2+\lvert\mathcal Q_{\rm LG}\rvert^2\)
from the horizon to infinity, for \(\mathcal Q_{\rm LG}(2)=1\) \\
\(F^{\rm H}_{\rm LG}\), \(a^{\rm H}_{\rm LG}\) & the coefficients of the
small-mass formulas of the energy and of the decay rate of a quasi-bound
state (E13) \\
\(y_{\rm LG}\) & in E13 to E19: an abbreviation of
\(\mathcal M_{\rm LG}m_{\rm LL}\) (in E15 to E19 with the mass at the
time considered) \\
\(S_{\rm H}\), \(H_{\rm LG}\) & the constant \(4\times4\) matrix
\(\mathrm{diag}(H_{\rm LG},H_{\rm LG})\) of the separated spinor,
\(H_{\rm LG}=i(\sigma_1+\sigma_3)/\sqrt2\) (E7, E15) \\
\(\sigma_1\), \(\sigma_3\) & two of the Pauli matrices \\
\(v\) & the advanced time of the Vaidya space-time (D10), constant along
every radial light ray that falls inward \\
\(\mathcal M_{\rm LG}(v)\) & in stage 6: the mass of the black hole as a
function of \(v\), in units of length \\
\(\mathcal P^{\rm V}_{\rm LG}(v,r)\),
\(\mathcal Q^{\rm V}_{\rm LG}(v,r)\) & the two radial functions of the
brochon in the coframe of D10 \\
\(p^{\rm V}_{\rm LG}(r)\), \(q^{\rm V}_{\rm LG}(r)\) & their profiles
for a constant mass,
\(\mathcal P^{\rm V}_{\rm LG}=e^{-i\omega^{\rm S}_{\rm LG}v}p^{\rm V}_{\rm LG}\),
\(\mathcal Q^{\rm V}_{\rm LG}=e^{-i\omega^{\rm S}_{\rm LG}v}q^{\rm V}_{\rm LG}\) \\
\(\Delta_{\rm LG}(r)\) & for a constant mass, \(v-\tau_{\rm LG}\):
\(\Delta_{\rm LG}'=1/(1+w_{\rm LG})\) \\
\(\chi_{\rm LG}\) & the rapidity \(-\ln(1+w_{\rm LG})\) of the boost
that turns the freely falling coframe into that of D10 \\
\(N^{\rm V}_{\rm LG}(v)\) & the charge of the brochon on the slice
\(v=\text{const}\): the integral of
\(2\lvert\mathcal P^{\rm V}_{\rm LG}\rvert^2\) over \(r\) \\
\(\delta r_{\rm LG}\) & the spacing of the grid of radii in stage 6 \\
\(\langle a,b\rangle_{\rm LG}\) & the charge product
\(\int2\bar a\,b\,dr\) of two radial functions on a slice \\
\(O_{n,{\rm LG}}\) & the occupation (projection) of the field on the
state \(n\) (E18) \\
\(p_n\) & in E18: the profile of the state \(nS_{1/2}\) of the final
mass on the slice \\
\(c_{n,{\rm LG}}\), \(W_{n,{\rm LG}}\) & the coefficient of the state
\(n\) in the decomposition of the field, and the charge
\(\lvert c_{n,{\rm LG}}\rvert^2\langle p_n,p_n\rangle_{\rm LG}\) that it
would carry alone \\
\(r_*\) & the tortoise coordinate
\(r+2\mathcal M_{\rm LG}\ln(r/2\mathcal M_{\rm LG}-1)\) \\
\(Z\), \(Z'\) & in E17: the nuclear charges of a hydrogen atom before
and after a sudden change \\
\(\alpha^{\rm Z}_{\rm LG}\) & the coupling
\((g^{\rm v}_{\rm LL})^2/(16\pi)\) of the zeugon between two
quasi-particles of opposite charge (D11) \\
\(\Lambda_{\rm LL}\) & the momentum cutoff of the earlier hypotheses,
the unit of energy of stage 7 \\
\(m^{\rm pair}_{\rm LG}\), \(\mu^{\rm pair}_{\rm LG}\) & the mass of one
particle of the pair of D11 (\(E^{\rm CJT,gap}_{\rm SF}\) or
\(m_{\rm LL}\)) and the reduced mass \(m^{\rm pair}_{\rm LG}/2\) \\
\(k^{\rm tid}_{\rm LG}\) & the strength
\(\mu^{\rm pair}_{\rm LG}\mathcal M_{\rm LG}/r^3\) of the tidal
potential on the relative coordinate of the pair (E20) \\
\(r^{\rm ion}_{\rm LG}\) & the radius at which the tidal field frees the
pair (classical estimate, E20) \\
\(M_{\rm Pl}\) & the Planck mass \((\hbar c/G_{\rm N})^{1/2}\) \\
\end{longtable}

\end{document}
